% physrevD.tex -- Physical Review D submission
\documentclass[aps,prd,twocolumn,showpacs,preprintnumbers,superscriptaddress,floatfix]{revtex4-2}

\usepackage{amsmath,amssymb,amsfonts,amsthm}
\usepackage{graphicx}
\usepackage{hyperref}
\usepackage{booktabs}

\newtheorem{theorem}{Theorem}
\newtheorem{proposition}[theorem]{Proposition}
\newtheorem{corollary}[theorem]{Corollary}
\newtheorem{lemma}[theorem]{Lemma}

\newcommand{\SU}{\mathrm{SU}}
\newcommand{\UU}{\mathrm{U}}
\newcommand{\SO}{\mathrm{SO}}
\newcommand{\Sp}{\mathrm{Sp}}
\newcommand{\Nf}{N_{\!f}}
\newcommand{\Nc}{N_{\!c}}
\newcommand{\Tr}{\mathrm{Tr}}
\newcommand{\adj}{\mathrm{adj}}
\newcommand{\diag}{\mathrm{diag}}
\newcommand{\keV}{\,\mathrm{keV}}
\newcommand{\MeV}{\,\mathrm{MeV}}
\newcommand{\GeV}{\,\mathrm{GeV}}
\newcommand{\TeV}{\,\mathrm{TeV}}
\newcommand{\ie}{i.e.\@}
\newcommand{\eg}{e.g.\@}
\newcommand{\zch}{\mathfrak{z}}

\begin{document}

\title{Composite fermions from $\SU(3)_F$ confinement:\\
  the sBootstrap as a structural effective field theory}

\author{A.~Rivero}
\email{rivero@unizar.es}
\affiliation{Departamento de F\'{\i}sica Te\'orica, Universidad de Zaragoza, 50009 Zaragoza, Spain}

\date{\today}

\begin{abstract}
We identify Standard Model fermions with composites of a
confining $\SU(3)_F$ family gauge theory: mesinos
($\psi_M$ of $M^i{}_j = Q^i\tilde{Q}_j$) are leptons;
adjugate composites provide quark partners via partial
compositeness.
The charge-squared mass law $m_i = \mu\,\zch_i^2$ is the
unique monomial for which $K = \sum\zch_i^2/(\sum|\zch_i|)^2$
is independent of the Cartan direction (Theorem~1);
the Seiberg seesaw then produces both a direct branch
($K = 2/3$ for leptons, $0.001\%$) and an inverse branch
($K^{-1} = 2/3$ for down quarks, $0.22\%$) from a single UV
condition.
An O'Raifeartaigh seed at $gv/m = \sqrt{3}$ breaks SUSY with
$K = 2/3$ exactly.
Five criteria---two structural, three suggestive---select
$N = 3$ generations.
% REVISED: REF-005
The principal open problem is the non-perturbative K\"ahler
metric, which controls the mapping from holomorphic to physical
masses.
\end{abstract}

\maketitle


%======================================================================
\section{Introduction}
\label{sec:intro}
%======================================================================

The Standard Model contains 13~free parameters in the flavour
sector---six quark masses, three charged-lepton masses, three CKM
angles, and a CP phase---with no known relations among them.
While the gauge sector is tightly constrained by the symmetry
group $\SU(3)_c \times \SU(2)_L \times \UU(1)_Y$, the Yukawa
couplings that determine fermion masses are arbitrary complex
numbers.
The absence of flavour structure is one of the outstanding
puzzles in particle physics, closely related to the hierarchy
problem (why $m_e/m_t \sim 3 \times 10^{-6}$?) and to the
origin of CP violation in the quark sector.

This paper develops a programme that addresses the flavour
problem not by postulating textures or discrete symmetries on the
Yukawa matrices, but by \emph{identifying the observed fermions
with composite fermions} of a confining gauge theory.
The key idea is that the mass hierarchies and mixing patterns are
not free parameters but consequences of the non-perturbative
dynamics of confinement: the Seiberg seesaw, the adjugate map,
and the charge algebra of $\UU(3)_F$ together constrain the
mass spectrum.

The idea that hadrons are composites of more fundamental quarks
was developed through the ``bootstrap'' programme of
Chew~\cite{Chew:1962}, which treated all hadrons
democratically as bound states of one another.
In 1966, Miyazawa~\cite{Miyazawa:1966} observed that the
spectrum of hadrons exhibits a \emph{supersymmetric} pattern,
with mesons ($q\bar{q}$, bosons) and baryons ($qqq$, fermions)
forming approximate supermultiplets.
Catto and G\"ursey~\cite{CattoGursey:1985} sharpened this: the
diquark $[qq]$ (in the antisymmetric colour
$\overline{\mathbf{3}}$) transforms as the SUSY partner of the
antiquark $\bar{q}$, so the baryon $q[qq]$ is literally the
fermionic partner of the meson $q\bar{q}$.

Masiero and Veneziano~\cite{MasieroVeneziano:1985} constructed
the first prototype model for one leptonic family from
$\SU(3)$ SQCD with $\Nf = \Nc = 3$, showing that soft SUSY
breaking splits composite supermultiplets to favour the
fermions; Karasik~\cite{Karasik:2020} revisited composite Higgs
from confining SQCD with modern tools.
% REVISED: REF-009
Luty and Mohapatra~\cite{LutyMohapatra:1997} constructed
explicit SUSY composite models where all quarks and leptons
emerge from confining $\SU(N)$ dynamics, with quarks as mesons
and leptons as baryons;
Luty and Terning~\cite{LutyTerning:2000} analysed the
K\"ahler metric problem for composite masses in
$\Nf = \Nc$ theories.
% REVISED: REF-009
Neither model has mass predictions; the present framework
differs in three respects: (i)~the confining group is
$\SU(3)_F$ with $\Nf = \Nc = 3$ (quantum-modified constraint,
not $\Nf = \Nc + 1$ dynamical superpotential); (ii)~the
identification is \emph{reversed} (mesinos = leptons, not
baryons); and (iii)~the charge-squared law provides
quantitative mass predictions.

The \emph{sBootstrap} (supersymmetric
bootstrap)~\cite{Rivero:2024epjc} takes this one step further.
Rather than treating the meson--baryon supersymmetry as an
approximate hadronic symmetry, it elevates it to a
\textbf{structural identification}: the composite fermions of a
confining gauge theory \emph{are} the Standard Model fermions.
The foundational statement is:
\begin{equation}
  \boxed{
  \text{mesinos} = \text{leptons}\,,
  \qquad
  \text{diquarkinos} = \text{quarks}\,.}
  \label{eq:sbootstrap_def}
\end{equation}
Here ``mesino'' means the fermionic component $\psi_M$ of the
meson chiral superfield $M^i{}_j = Q^i\tilde{Q}_j$, where
$Q^i$ and $\tilde{Q}_j$ are the fundamental ``preon'' fields of
a confining $\SU(3)_F$ family gauge theory.
The ``diquarkino'' is the fermionic component of the adjugate
composite $\adj(M)_i$, whose quantum numbers match those of a
diquark-type operator.

\paragraph{Evidence: $m_\pi^2 - m_\mu^2 \approx f_\pi^2$.}
The mass relation $m_\pi^2 - m_\mu^2 \approx f_\pi^2$
is \emph{evidence} for the identification, not its definition.
Numerically, $m_\pi^2 - m_\mu^2 = 139.57^2 - 105.66^2 =
8316\MeV^2$ while $f_\pi^2 = 92.2^2 = 8501\MeV^2$ ($2.2\%$
discrepancy).
In the Volkov--Akulov framework~\cite{VolkovAkulov:1973},
spontaneous SUSY breaking gives the mass relation
$m_{\rm fermion}^2 = m_{\rm boson}^2 - f^2$, where $f$ is the
SUSY-breaking scale.
If we identify $f = f_\pi$, then the muon is a pseudo-Goldstino
whose mass is split from the pion by the chiral symmetry
breaking scale.

\paragraph{The emergent $\mathcal{N}{=}2$ pairing.}
In $\mathcal{N}{=}1$ SUSY, each chiral multiplet
$(M_i, \psi_{M_i})$ provides one fermion (the mesino).
But the confined spectrum also contains the adjugate composite
$\adj(M)_i$ (Sec.~\ref{sec:seiberg}), which forms a second
chiral multiplet per family.
The pair $(M_i, \adj(M)_i)$ has the mathematical structure of an
$\mathcal{N}{=}2$ hypermultiplet---not because the UV theory has
$\mathcal{N}{=}2$ supersymmetry, but because confinement
\emph{automatically} produces both $M$ and $\adj(M)$ with the
correct quantum numbers.

\paragraph{The Q-generator problem.}
In the Miyazawa superalgebra, a single supercharge $Q$ maps
$q\bar{q} \to q[qq]$ (meson to baryon).
In the sBootstrap, the involution $M_i \leftrightarrow
\Lambda^6/M_i$ (the adjugate map) is an \emph{exact}
$\mathbb{Z}_2$ symmetry of the quantum-modified constraint
surface $\det M = \Lambda^6$: it exchanges the direct and
inverse branches and is a holomorphic identity, not an
approximation.
The \textbf{open question} is whether this $\mathbb{Z}_2$ can
be promoted to a continuous $\SU(2)_R$ rotation within the
doublet $(M_i, \adj(M)_i)$.
If so, the associated generator $Q_2$ would connect the meson
sector (leptons) to the adjugate sector (quarks), realising
Miyazawa's vision with two supercharges.
$Q_1$ is the standard $\mathcal{N}{=}1$ supercharge; the
approximate $\SU(2)_R$ is broken by the adjoint mass $\mu_\Phi$.
Constructing $Q_1$ and $Q_2$ explicitly in the confined theory
is an open problem.

\paragraph{Summary of results.}
The framework makes the following quantitative connections.
The energy-balance condition $K = 2/3$ is satisfied by
charged leptons ($0.001\%$), pseudoscalar mesons $(-\pi, D_s, B)$
($0.10\%$), inverse down quarks ($0.22\%$), and the quark chain
$(t,b,c)$--$(b,c,-s)$--$(c,s,0)$--$(s,0,m_u{+}m_d)$ with joint
significance $p \sim 1/21\,000$.
The bion relation $\zch_b = 3\zch_s + \zch_c$ predicts
$m_b = 4183\MeV$ ($0.01\%$), and the top completion formula
predicts $m_t = 168.5\GeV$ ($2.3\%$).
Five criteria (two structural, three suggestive) select
$N = 3$ generations.
The framework does NOT yet predict: the CKM matrix beyond the
Cabibbo angle, the PMNS matrix, the origin of the Cartan angle
$\alpha$, or the K\"ahler metric $Z_i$.

\paragraph{Plan of the paper.}
We proceed by building up the framework layer by layer:
composites (\S\ref{sec:composites}) $\to$
Seiberg dynamics (\S\ref{sec:seiberg}) $\to$
charge-squared law (\S\ref{sec:charges}) $\to$
emergent isospin (\S\ref{sec:N2}) $\to$
SUSY breaking (\S\ref{sec:susy_breaking}) $\to$
inverse sector and chain (\S\ref{sec:inverse}) $\to$
generation uniqueness (\S\ref{sec:SO32}) $\to$
Lagrangian (\S\ref{sec:lagrangian}) $\to$
discussion and conclusions (\S\ref{sec:discussion},
\S\ref{sec:open}).


%======================================================================
\section{Composites and confinement}
\label{sec:composites}
%======================================================================

\subsection{The gauge-theoretic realisation}

The UV gauge group is
\begin{equation}
  G_{\rm UV} = \SU(3)_F \times \SU(3)_c \times \SU(2)_L
  \times \UU(1)_Y\,,
  \label{eq:GUV}
\end{equation}
where $\SU(3)_F$ is the \emph{family gauge symmetry} that
confines at a scale $\Lambda_F$.
The term ``family'' refers to the three generations of SM
fermions: $\SU(3)_F$ acts on the generation index
$i \in \{1, 2, 3\}$.

The fundamental matter fields (``preons'') are:
\begin{itemize}
  \item $Q^\alpha_i$: transforming as
    $(\mathbf{3}_F, \mathbf{1}_c, \mathbf{3}_L)$,
    where $\alpha \in 3_L$ is an index under the
    \emph{global} $\SU(3)_L$ flavour symmetry
    (containing $\SU(2)_L \times \UU(1)_Y$ as a subgroup).
  \item $\tilde{Q}_i^\alpha$: the conjugate,
    $(\overline{\mathbf{3}}_F, \mathbf{1}_c,
    \overline{\mathbf{3}}_L)$.
\end{itemize}
The $\SU(3)_L$ is \textbf{not} a gauge symmetry: it is a global
symmetry whose sole role is to organise the branching
$\mathbf{3}_L \to \mathbf{2}_{-1/2} \oplus \mathbf{1}_{+1}$,
which assigns the correct SM quantum numbers to the composites
and projects out exotic charges
(Sec.~\ref{sec:SO32}).
The preons are colour singlets: quarks (which carry colour) arise
through partial compositeness.

\begin{table}
\centering\small
\caption{UV field content (above $\Lambda_F$).
  All fields listed as $\mathcal{N}{=}1$ chiral superfields
  except the gauge multiplets.}
\label{tab:UV_fields}
\begin{tabular}{lccccl}
\hline
Field & $\SU(3)_F$ & $\SU(3)_c$ & $\SU(2)_L$ & $\UU(1)_Y$ & Role \\
\hline
$Q^\alpha_i$ & $\mathbf{3}$ & $\mathbf{1}$ & $\mathbf{2}\oplus\mathbf{1}$ & $-\tfrac{1}{2},+1$ & preon \\
$\tilde{Q}_i^\alpha$ & $\overline{\mathbf{3}}$ & $\mathbf{1}$ & $\overline{\mathbf{2}}\oplus\mathbf{1}$ & $+\tfrac{1}{2},-1$ & anti-preon \\
$\Phi$ & $\mathbf{8}$ & $\mathbf{1}$ & $\mathbf{1}$ & $0$ & adjoint (aux.) \\
$V_F^\mu$ & $\mathbf{8}$ & $\mathbf{1}$ & $\mathbf{1}$ & $0$ & gauge boson \\
$q_L^a$ & $\mathbf{1}$ & $\mathbf{3}$ & $\mathbf{2}$ & $+\tfrac{1}{6}$ & elem.\ quark \\
$u_R^a$ & $\mathbf{1}$ & $\mathbf{3}$ & $\mathbf{1}$ & $+\tfrac{2}{3}$ & elem.\ quark \\
$d_R^a$ & $\mathbf{1}$ & $\mathbf{3}$ & $\mathbf{1}$ & $-\tfrac{1}{3}$ & elem.\ quark \\
\hline
\end{tabular}
\end{table}


\subsection{The meson decomposition}

When $\SU(3)_F$ confines, the preons form composite bound states.
The meson superfield $M^\alpha{}_i = Q^\alpha\tilde{Q}_i$ carries
one $\SU(3)_L$ index and one $\SU(3)_F$ index, giving
$3 \times 3 = 9$ chiral multiplets.
Under the electroweak branching
$\mathbf{3}_L \to \mathbf{2}_{-1/2} \oplus \mathbf{1}_{+1}$:
\begin{equation}
  M^\alpha{}_i \;=\;
  \underbrace{M^{1,2}{}_i}_{(\mathbf{2}_{-1/2})}
  \;\oplus\;
  \underbrace{M^{3}{}_i}_{(\mathbf{1}_{+1})}\,.
  \label{eq:meson_decomp}
\end{equation}
This decomposition is a \textbf{projection tool}: its role is
to clean the $\mathbf{15} \oplus \overline{\mathbf{15}} \oplus
\mathbf{24}$ composite content for the $\SO(32)$ string embedding
(Sec.~\ref{sec:SO32}).
Without it, the $\mathbf{24}$ of $\SU(5)_f$ contains
$(\overline{\mathbf{3}}, \mathbf{2})_{5/6}$ with exotic charge
$Q = 4/3$~\cite{Rivero:2024epjc}---the analogue of $X$ and $Y$
bosons in $\SU(5)$ GUT.
The branching projects these out, restricting composites to
SM-compatible charges.
This is the \emph{only} reason for introducing $\SU(3)_L$.
The mesino identification follows:
\begin{equation}
  \boxed{
  \psi_{M^{1,2}{}_i} \equiv L_i
  \;=\; \begin{pmatrix} \nu_i \\ e_i \end{pmatrix}_L\,,
  \qquad
  \psi_{M^{3}{}_i} \equiv e^c_i\,.}
  \label{eq:mesinos_leptons}
\end{equation}
The 9~mesino fermions furnish the complete SM lepton sector:
$3L_i + 3e^c_i$ $=$ 3 lepton doublets $+$ 3 charged-lepton
singlets.
Anomaly matching between UV preons and IR composites is satisfied
for $\Nf = \Nc$~\cite{CsakiSchmaltzSkiba:1997}.


\begin{table}
\centering\small
\caption{IR composites (below $\Lambda_F$, $\Nf = \Nc = 3$).
  Independent DOF: $9\,(M) + 1\,(B) + 1\,(\tilde{B})
  - 1\,(X) = 10 = 18 - 8$.
  $\adj(M)$ is \emph{derived}---not independent DOF.}
\label{tab:IR_composites}
\begin{tabular}{llcccl}
\hline
Composite & Expr.\ & $\SU(2)_L$ & $\UU(1)_Y$ & Ch.\ & SM id.\ \\
\hline
$M^{1,2}{}_i$ & $Q^{1,2}\tilde{Q}_i$ & $\mathbf{2}$ & $-\tfrac{1}{2}$ & 6 & $H_d$, $L_i$ \\
$M^{3}{}_i$ & $Q^{3}\tilde{Q}_i$ & $\mathbf{1}$ & $+1$ & 3 & $e^c_i$ \\
\hline
\multicolumn{6}{l}{\textit{Derived operators:}} \\
$\adj(M)^{1,2}{}_i$ & $\varepsilon\varepsilon MM$ & $\mathbf{2}$ & $+\tfrac{1}{2}$ & (6) & $H_u$ \\
$\adj(M)^{3}{}_i$ & $\varepsilon\varepsilon MM$ & $\mathbf{1}$ & $-1$ & (3) & $\tilde\nu$-type \\
$B$ & $\varepsilon QQQ$ & $\mathbf{1}$ & $0$ & 1 & sterile \\
$\tilde{B}$ & $\varepsilon\tilde{Q}^3$ & $\mathbf{1}$ & $0$ & 1 & sterile \\
$X$ & constraint & $\mathbf{1}$ & $0$ & $-1$ & Lagr.\ mult.\ \\
\hline
\end{tabular}
\end{table}


\subsection{Quarks via partial compositeness}
\label{sec:composites_quarks}

Quarks carry colour ($\SU(3)_c$) and cannot be mesinos (which are
colour singlets).
% REVISED: REF-006 — replaced Fradkin-Shenker with SUSY complementarity
In supersymmetric gauge theories, the smooth connection between
confinement and Higgs phases---the SUSY analogue of
Fradkin--Shenker complementarity~\cite{FradkinShenker:1979}---is
established by Seiberg
duality~\cite{Seiberg:1994bz} and s-confinement: for $\SU(3)$
with $\Nf = 3$, the confined composites ($M, B, \tilde{B}$,
subject to the quantum constraint) describe the same 10 light DOF
as the Higgs phase where preons develop VEVs and eat 8 Goldstone
bosons ($2 \times 9 - 8 = 10$).

% REVISED: REF-006
\paragraph{Colour-$\mathbf{3}$ impossibility.}
Since the preons $Q^i$ and $\tilde{Q}_i$ are colour singlets
(Table~\ref{tab:UV_fields}), every polynomial composite
$Q^{i_1}\cdots Q^{i_n}\tilde{Q}_{j_1}\cdots\tilde{Q}_{j_m}$ is
a product of colour singlets and hence a colour singlet.
In particular, $\SU(3)_F$ composites \emph{never} carry colour
$\mathbf{3}$---only $\mathbf{1}$, $\mathbf{6}$, $\mathbf{8}$,
etc.\ appear.
Quarks, which carry colour, therefore \emph{cannot} be composites
of this theory.

\paragraph{Quarks via partial compositeness.}
The coloured quarks of the SM enter through \textbf{partial
compositeness}~\cite{Kaplan:1991dc}:
each elementary quark $q_{\rm elem}$ (which carries colour and
is a singlet of $\SU(3)_F$, see Table~\ref{tab:UV_fields})
mixes linearly with a composite operator $\mathcal{O}$ of
matching SM quantum numbers through both left- and right-handed
couplings:
\begin{equation}
  \mathcal{L}_{\rm mix} = \lambda_{q_L}\,
  \bar{q}_{L,\rm elem}\,\mathcal{O}_L
  + \lambda_{q_R}\,
  \bar{q}_{R,\rm elem}\,\mathcal{O}_R
  + \text{h.c.}
  \label{eq:partial_comp}
\end{equation}
The physical quark mass is~\cite{Kaplan:1991dc}
\begin{equation}
  m_q \sim \frac{\lambda_{q_L}\,\lambda_{q_R}}{M_*}
  \,\langle\mathcal{O}\rangle\,,
  \label{eq:mq_partial}
\end{equation}
where $M_*$ is the compositeness scale.
If both mixing parameters scale with the family charge,
$\lambda_{q_L} \propto \lambda_{q_R} \propto \zch_q$, then
$m_q \propto \zch_q^2$---the charge-squared law.
The simplest realisation has $\lambda_{q_L} = \lambda_{q_R}$
(parity-symmetric mixing), but even with unequal mixings the
$\zch^2$ scaling holds provided both are proportional to the
same family charge.
By Seiberg duality~\cite{Seiberg:1994bz,CsakiSchmaltzSkiba:1997},
the confinement and Higgs phases of $\SU(3)_F$ are smoothly
connected (complementarity), so the physical quark in the
confined description IS the dressed preon in the Higgs
description.

\paragraph{The lepton sector obstruction.}
There is an honest open problem.
The holomorphic meson masses from the Seiberg vacuum
(Sec.~\ref{sec:seiberg}) give mass \emph{ratios}
determined by the UV preon masses.
The mesino mass ratio
$m_\psi(M^u{}_s)/m_\psi(M^u{}_d) = m_s/m_d \approx 20$,
while the physical lepton ratio $m_\mu/m_e = 207$---a factor
$10.3$ discrepancy.
K\"ahler corrections preserve the rank of the mass matrix but
can only introduce flavour-independent rescalings if $Z$ is
generation-universal.
Flavour-\emph{dependent} rescaling (different $Z_i$ per
generation) can generate the discrepancy.
Two candidate mechanisms:
(i) Nelson--Strassler anomalous dimensions at a confining
fixed point~\cite{NelsonStrassler:2000};
(ii) a separate $\SU(2)$ confining sector with $\Nf = 3$ and
an O'Raifeartaigh deformation (Sec.~\ref{sec:susy_breaking}).
We flag this as an open issue, not an inconsistency: the
framework predicts the \emph{structure} of the mass matrix
(diagonal, with entries $\propto \zch_i^2$) but not the
overall K\"ahler normalisation.


%======================================================================
\section{Seiberg confinement and the adjugate mechanism}
\label{sec:seiberg}
%======================================================================

\subsection{Review: s-confinement for $\Nf = \Nc$}

We review the non-perturbative dynamics of $\SU(\Nc)$ SQCD with
$\Nf$ flavours of quarks and antiquarks in the $\mathcal{N}{=}1$
framework~\cite{Seiberg:1994bz,Tong:SUSY}.

For $\Nf = \Nc$, the theory is \emph{s-confining} (``smoothly
confining''): the low-energy degrees of freedom are the
gauge-invariant composites $M$, $B$, $\tilde{B}$, and there is
no residual gauge group.
The classical moduli space is parametrised by the constraint
$\det M - B\tilde{B} = 0$.
Quantum effects (a one-instanton calculation) modify this to the
\textbf{quantum-modified constraint}:
\begin{equation}
  \det M - B\tilde{B} = \Lambda^{2\Nc}\,,
  \label{eq:constraint}
\end{equation}
where $\Lambda$ is the dynamical scale.
The right-hand side $\Lambda^{2\Nc}$ is generated by instantons
and is exact (protected by holomorphy and symmetries).
The effect is to smooth out the moduli space: the origin
$M = B = \tilde{B} = 0$ is removed.

The constraint is implemented in the effective superpotential by
a Lagrange multiplier $X$:
\begin{equation}
  W = X\bigl(\det M - B\tilde{B} - \Lambda^6\bigr)
    + \sum_i m_i\,M^i{}_i\,.
  \label{eq:Wtot}
\end{equation}
The first term enforces the constraint; the second adds explicit
preon masses which break global symmetry and lift the moduli
space to isolated vacua.


\subsection{The mesonic vacuum and the Seiberg seesaw}

On the mesonic branch ($B = \tilde{B} = 0$), we solve the
$F$-term equations $\partial W / \partial M^i{}_j = 0$.
For diagonal $M = \diag(M_1, M_2, M_3)$:
\begin{equation}
  \frac{\partial W}{\partial M_i}
  = X \cdot \frac{\partial(\det M)}{\partial M_i} + m_i
  = X \cdot \prod_{j \neq i} M_j + m_i = 0\,.
  \label{eq:Fterm_M}
\end{equation}
The product $\prod_{j \neq i} M_j$ is precisely the $i$-th
diagonal entry of the adjugate matrix.
From the constraint $\det M = M_1 M_2 M_3 = \Lambda^6$,
and taking ratios of $F$-term equations for different $i$:
$M_i/M_j = m_j/m_i$, which gives the \textbf{Seiberg seesaw}:
\begin{equation}
  \langle M_i \rangle = \frac{\Lambda^2 P^{1/3}}{m_i}\,,
  \qquad P = m_1 m_2 m_3\,.
  \label{eq:vacuum}
\end{equation}
Heavy preons produce light mesons, and vice versa.
This is the single most important equation in the framework.


\subsection{The $\Nf$-dependent phase structure}

The behaviour of $\SU(3)_F$ depends critically on how many preon
flavours are light compared to the RG scale.
As the energy decreases, heavy quarks are sequentially integrated
out, changing the effective number of flavours and therefore the
phase of the gauge theory.

\begin{table}
\centering\small
\caption{Phase structure of $\SU(3)_F$ as heavy flavours decouple.
  Scale matching: $\Lambda_3^6 = \Lambda_4^5 m_c$,
  $\Lambda_4^5 = \Lambda_5^4 m_b$.
  Each threshold crossing changes the non-perturbative dynamics
  qualitatively.}
\label{tab:phases}
\begin{tabular}{clll}
\hline
$\Nf$ & Energy range & Phase & Key property \\
\hline
$6$ & $Q > m_t$ & conformal window & IR fixed point \\
$5$ & $m_b < Q < m_t$ & conformal window & IR fixed point \\
  &  & ($4.5 < 5 < 9$) & \\
$4$ & $m_c < Q < m_b$ & ISS metastable~\cite{ISS:2006} & unique integer \\
  &  & & in $(\Nc, \tfrac{3}{2}\Nc)$ \\
$3$ & $Q < m_c$ & s-confinement & $\det M = \Lambda^6$ \\
\hline
\end{tabular}
\end{table}

At $\Nf = 5$, the theory enters the conformal window
($\frac{3}{2}\Nc < \Nf < 3\Nc$, i.e.\ $4.5 < 5 < 9$): it flows
to an interacting IR fixed point, the Seiberg conformal fixed
point.
At $\Nf = 4$, it saturates the lower bound of the ISS
window~\cite{ISS:2006}: $\Nc {+} 1 \leq \Nf < \frac{3}{2}\Nc$,
i.e.\ $4 \leq \Nf < 4.5$.
% REVISED: REF-010
This is the window where metastable SUSY-breaking vacua exist.
$\Nf = 4$ is the \emph{unique} integer in this window for
$\Nc = 3$---a discrete selection that plays a role in the
effective Lagrangian (Sec.~\ref{sec:lagrangian}).
The metastable vacuum has bounce action
$S_{\rm bounce} \sim 10^2$--$10^3$, giving a lifetime vastly
exceeding the age of the universe.
At $\Nf = 3$, the theory s-confines with the quantum-modified
constraint described above.


\subsection{The adjugate: definition and properties}

The adjugate of a square matrix $M$ is defined by:
\begin{equation}
  \adj(M) \cdot M = \det(M) \cdot I\,.
  \label{eq:adj_def}
\end{equation}
Unlike the inverse $M^{-1} = \adj(M)/\det M$, the adjugate
is a \emph{polynomial} in the entries of $M$---holomorphic and
well-defined even when $\det M = 0$.
For diagonal $M = \diag(M_1, M_2, M_3)$:
\begin{equation}
  \adj(M) = \diag(M_2 M_3,\; M_1 M_3,\; M_1 M_2)\,.
  \label{eq:adj_3x3}
\end{equation}
Each entry is the product of the \emph{other two} diagonal
entries.
Verification: $\adj(M)_1 \cdot M_1 = M_1 M_2 M_3 = \det M$.
The ordering is \emph{reversed}: if
$M_1 < M_2 < M_3$, then
$\adj(M)_1 = M_2 M_3 > \adj(M)_2 = M_1 M_3 > \adj(M)_3 =
M_1 M_2$.
This ordering reversal is the physical origin of the
direct/inverse mass duality: the heaviest preon produces the
lightest meson VEV and the largest adjugate entry.
For general matrices, $\adj(M)_{ij} = (-1)^{i+j} C_{ji}$
(transpose of the cofactor matrix).
The key mathematical property is that $\adj(M)$ is holomorphic:
it is well-defined even at $\det M = 0$, unlike $M^{-1}$.


\subsection{The two Yukawa operators: direct vs.\ inverse}

From the $F$-term equations~\eqref{eq:Fterm_M}:
\begin{equation}
  \adj(M)_i = -\frac{m_i}{X} = -\Lambda^3\,m_i\,,
  \label{eq:adjM_from_Fterm}
\end{equation}
using $X = -P^{1/3}/\Lambda^4$ from the constraint
(the simpler form $X = -1/\Lambda^3$ holds in the convention
$P^{1/3} = \Lambda$, \ie\ equal UV masses).
Together with $\langle M_i \rangle = \Lambda^2 P^{1/3}/m_i$:
\begin{align}
  M_i &\propto \frac{1}{m_i}
    \quad\text{(inverse: heavy UV $\to$ light IR)}\,,
    \label{eq:M_inverse}\\
  \adj(M)_i &\propto m_i
    \quad\text{(direct: heavy UV $\to$ heavy IR)}\,.
    \label{eq:adjM_direct}
\end{align}
These are the two natural Yukawa operators of the confined
theory.
The Seiberg dynamics provides both from the single composite
field $M$, related by the holomorphic identity
$\adj(M) = \Lambda^6 M^{-1}$ on the constraint surface.

The physical assignment~\cite{Seiberg:1994bz,CsakiSchmaltzSkiba:1997}:
\begin{itemize}
  \item \textbf{Leptons} are mesinos $\psi_M$.
    Their mass comes from the $F$-term
    $|F_{M_i}| \propto m_i$ (direct, proportional to UV mass).
    Hence: $K(m_e, m_\mu, m_\tau) = 0.666661$ (standard
    energy-balance, $0.001\%$).
  \item \textbf{Down-type quarks} couple to the meson VEV
    $\langle M \rangle \propto 1/m_i$ through partial
    compositeness.
    Hence: $K^{-1}(m_d, m_s, m_b) = 0.6652$ (inverse
    energy-balance, $0.22\%$).
\end{itemize}
A \emph{single UV condition} $K(m^{\rm UV}) = 2/3$ produces
both the standard lepton Koide and the inverse quark Koide,
unified by the involution $M_i \leftrightarrow \Lambda^4/M_i$.

\paragraph{Automatic pairing of direct and inverse.}
If the UV masses satisfy $m_i = \mu\zch_i^2$ with $\UU(3)_F$
charges, then on the constraint surface
$a_i = \alpha\zch_i^2$ and $d_i = \beta/\zch_i^2$ with
$\alpha\beta = \Lambda^6$.
Both $K_{\rm dir} = K[\zch_1,\zch_2,\zch_3]$ and
$K_{\rm inv} = K[\zch_1^{-1},\zch_2^{-1},\zch_3^{-1}]$
evaluate to $2/3$ by the same algebraic identity
(Theorem~\ref{thm:n2}): the $n = 2$ uniqueness holds for
\emph{both} $\zch_i$ and $1/\zch_i$.
The inverse energy-balance is a holomorphic consequence of the
quantum-modified vacuum, not an independent phenomenological
input.

\begin{table}
\centering\small
\caption{Direct and inverse branch assignment.
  Product rule: $a_i \cdot d_i = \mu\nu = \text{const}$
  ($\det M = \Lambda^6$).}
\label{tab:branches}
\begin{tabular}{lllll}
\hline
Branch & Operator & Formula & Order & Sector \\
\hline
Direct & $\adj(M)_i$ & $a_i = \mu\zch_i^2$ & $a_1 < a_2 < a_3$ & Leptons \\
Inverse & $M_i$ & $d_i = \nu/\zch_i^2$ & $d_1 > d_2 > d_3$ & Down quarks \\
\hline
\end{tabular}
\end{table}


\subsection{Off-diagonal masses and transparency}

The off-diagonal meson masses
are~\cite{CsakiSchmaltzSkiba:1997}:
\begin{equation}
  m(M^i{}_j) = \frac{m_i m_j}{P^{1/3}}\,,
  \qquad i \neq j\,.
  \label{eq:offdiag}
\end{equation}
These depend only on ratios of the input masses---they are
$\Lambda$-independent.
The diagonal mass matrix has $(W_S)_{kk} = 0$ for all $k$
(by multilinearity of the determinant: expanding $\det M$ along
row $k$, the derivative $\partial(\det M)/\partial M_k^k$
contains no factor of $M_k^k$, so it cancels in the full
$F$-term).
This means the Seiberg dynamics is \textbf{transparent}: it
preserves $K = 2/3$ (if present in the UV masses) but cannot
generate it.
The origin of $K = 2/3$ must come from the charge-squared law
(Sec.~\ref{sec:charges}).


%======================================================================
\section{The charge-squared mass law}
\label{sec:charges}
%======================================================================

We derive the mass ansatz $m_i = \mu\,\zch_i^2$, where
$\zch_i = z_0 + z_i$ is the signed family charge.
Three independent arguments converge on the same result.

\subsection{The $\UU(3)_F$ charge algebra}
\label{sec:charge_algebra}

% REVISED: REF-011 — sign conventions stated up front
\paragraph{Sign conventions in the energy-balance formula.}
The Koide ratio $K = \sum m_i / (\sum\sqrt{m_i})^2$ is defined for
positive masses.
Throughout this paper, the ``signed charge'' $\zch_i$ can be
negative when the corresponding particle is a pseudo-Goldstone
boson approaching $m = 0$ from the chiral limit
(e.g.\ the pion), or when a mass triple contains both direct and
inverse-branch members (e.g.\ the quark chain link $(b, c, -s)$).
In these cases, the energy-balance formula is evaluated with
$m_i = \mu\,\zch_i^2$ and $\sqrt{m_i} = |\zch_i|
\cdot \mathrm{sign}(\zch_i)$, so that \emph{signed} square roots
enter the denominator.
The sign assignment is \emph{not} freely chosen to fit: it is
determined by the branch (direct vs.\ inverse, Sec.~\ref{sec:seiberg})
or by the chiral limit structure (GOR relation, $m_\pi^2 \propto
m_u + m_d \to 0$).

The family symmetry $\UU(3)_F = \SU(3)_F \times \UU(1)$ has a
Cartan subalgebra spanned by two diagonal generators of
$\SU(3)_F$ plus the $\UU(1)$ generator.
With canonical normalisation:
\begin{equation}
  \sum_{i=1}^{3} z_i = 0\,,\qquad
  \sum_{i=1}^{3} z_i^2 = \frac{1}{2}\,,\qquad
  z_0 = \frac{1}{\sqrt{6}}\,.
\end{equation}
The signed charge is $\zch_i = z_0 + z_i$, with
$z_i = \cos(\alpha + 2\pi(i-1)/3)/\sqrt{3}$ parametrised by
the Cartan direction $\alpha$ (the ``sector angle'').


\subsection{Argument 1: Algebraic uniqueness of $n = 2$}
\label{sec:uniqueness_n2}

Consider a general monomial mass law $m_i = c\,(z_0 + z_i)^n$
and define:
\begin{equation}
  K_n(\alpha) \equiv
  \frac{\sum_i (z_0{+}z_i)^n}
  {\bigl(\sum_i (z_0{+}z_i)^{n/2}\bigr)^2}\,.
  \label{eq:Kn}
\end{equation}

\begin{theorem}[Uniqueness of $n = 2$]
\label{thm:n2}
For integer powers $n$ and $\UU(3)$ charges $\zch_i = z_0 + z_i$,
the ratio $K_n(\alpha)$ is independent of the Cartan direction
$\alpha$ if and only if $n = 2$ (besides $n = 0$).
At $n = 2$: $K_2 = 2/3$.
\end{theorem}

\begin{proof}
We expand $\sum(z_0 + z_i)^n$ using the binomial theorem.
The power sums $S_k = \sum z_i^k$ satisfy:
$S_0 = 3$; $S_1 = 0$ (tracelessness of $\SU(3)$);
$S_2 = 1/2$ (normalisation, $\alpha$-independent);
and $S_k$ for $k \geq 3$: all depend on $\alpha$
through $\cos 3\alpha$ or higher harmonics.

For $n = 2$:
$\sum (z_0+z_i)^2 = 3z_0^2 + 2z_0 \cdot 0 + 1/2 = 1$.
Also $\sum (z_0+z_i) = 3z_0 = \sqrt{3/2}$.
So $K_2 = 1/(3/2) = 2/3$, independent of $\alpha$.

For $n \geq 3$: the binomial expansion of the numerator contains
the term $\binom{n}{3}z_0^{n-3}S_3$ with coefficient
$n(n-1)(n-2)/6 \neq 0$.
Since $S_3 = \sum z_i^3 \propto \cos 3\alpha$ (using
$z_i = \cos(\alpha + 2\pi(i-1)/3)/\sqrt{3}$), the numerator
acquires $\alpha$-dependence.
The denominator $(\sum\zch_i^{n/2})^2$ generically also depends
on $\alpha$ for $n \geq 3$, and the $\alpha$-dependence does not
cancel in the ratio.
For $n = 1$: $K_1 = \sum\zch_i/(\sum\zch_i^{1/2})^2$ involves
square roots of the charges, which are not polynomial in $\alpha$;
the ratio still varies.

Computer algebra (Sympy) verification confirms $K_n$ varies for
all tested real $n \in [-5, 10]$ at spacing $0.001$, except
$n = 2$ and the trivial $n = 0$.
\end{proof}

The physical meaning: $n = 2$ is the \emph{unique} monomial mass
law for which the energy-balance ratio $K$ is a group-theoretic
identity ($K = 2/3$), independent of which sector we are looking
at.
For any other exponent, $K$ would hold in some sectors but not
others---destroying the universality observed across leptons,
mesons, and quarks.

\subsection{Argument 2: Composite self-energy}

Koide's original derivation~\cite{Koide:1983qe} provides a
physical mechanism.
The Coulomb self-energy of a composite with total family charge
$\zch_i$ is:
\begin{equation}
  \delta_i = \frac{1}{2}\bigl[Q_F(A) + Q_F(B)\bigr]^2 \cdot v_F
  \propto \zch_i^2\,,
  \label{eq:koide_selfe}
\end{equation}
the gauge-theoretic version of the capacitor energy
$E = Q^2/(2C)$: Gauss's law gives $|\mathbf{E}| \propto \zch_i$,
and the energy $U \propto \int |\mathbf{E}|^2 \propto \zch_i^2$.
In a confining theory, the dominant mass contribution is the
string tension $\sigma \cdot L$, which is generation-independent
(it depends on the colour representation, not the Cartan
eigenvalue).
The $\zch_i^2$ scaling applies to the generation-dependent
Coulomb self-energy at short distances.
The generation-independent piece is absorbed into the overall
scale $\mu$.
In SUSY, the fermion mass is an $F$-term (holomorphic),
protected against non-perturbative confining corrections.


\subsection{Argument 3: Bilinear structure}

The meson $M^i{}_i = Q^i\tilde{Q}_i$ is a bilinear.
If each preon contributes $\propto \zch_i$, the composite VEV
scales as $\zch_i^2$~\cite{Koide:1990}.
This is the same structure as a seesaw: any mass from two VEV
insertions gives $m \propto v^2$.

\paragraph{Status summary.}
Argument~1 is a rigorous algebraic theorem.
Argument~2 is structural motivation (Coulomb scaling
assumed to survive confinement).
Argument~3 is a consistency check.
The charge-squared law is motivated by three structural inputs:
(i)~\textbf{Holomorphy}: the superpotential $W(M)$ is holomorphic,
so the mass is $\propto \zch^2$ (holomorphic square), not
$|\zch|^2$ (Hermitian norm);
(ii)~\textbf{Bilinear structure}: $M = Q\tilde{Q}$ involves two
preon fields each contributing $\propto \zch$;
(iii)~\textbf{Canonical normalisation}: $z_0 = 1/\sqrt{6}$
from the $\UU(3)_F$ kinetic term; the $n = 2$ uniqueness
(Theorem~\ref{thm:n2}) shows this is the only monomial giving
$K = 2/3$ universally.

The charge-squared law $m_i = \mu\,\zch_i^2$ is the framework's
central \emph{postulate}: uniquely motivated but not derived from
confining dynamics.
The uniqueness theorem is a \emph{necessary} condition: if masses
are monomials, then $n = 2$ is selected.
It does not exclude non-monomial $m = f(\zch)$.
Seven SUSY-breaking derivation attempts fail because gauge
mediation gives $\zch^2$ scaling for soft (K\"ahler) masses,
not for holomorphic (superpotential) masses needed for exact
$K = 2/3$.


\subsection{The energy-balance condition}

For three charges $\zch_k = z_0 + z_k$ with $\sum z_k = 0$:
\begin{equation}
  K = \frac{\sum \zch_k^2}{(\sum \zch_k)^2}
  = \frac{1}{3} + \frac{\langle z_k^2\rangle}{3z_0^2}\,.
  \label{eq:K_derivation}
\end{equation}
Therefore:
\begin{equation}
  \boxed{K = \frac{2}{3}
  \quad\Longleftrightarrow\quad
  \langle z_k^2\rangle = z_0^2\,.}
  \label{eq:energy_balance}
\end{equation}
At $K = 2/3$, the variance of the $\SU(3)$ charges equals the
square of the $\UU(1)$ charge: the splitting scale matches the
base scale.
This is the \emph{equipartition point} between the perturbative
limit ($K \to 1/3$, all masses equal) and the dominant limit
($K \to 1$, one mass dominates).

In the charge plane, the variable
$\langle z^2\rangle/z_0^2 = 3K - 1$ measures deviation from
$K = 2/3$.
The zero-mass algebraic fixed points of the charge plane occur at
$\alpha = 15^\circ$ and $45^\circ$ ($\SU(3)$ convention), with
ratio $r_0 = (2+\sqrt{3})^2 \approx 13.93$.
All five physical Koide tuples lie within $5.5^\circ$ of the
$15^\circ$ fixed point---their clustering reflects proximity to
a single algebraic attractor.

$K = 2/3$ is \textbf{not} an RG attractor: one-loop soft-mass
renormalisation drives $K \to 1/3$, not $K \to 2/3$.
The energy-balance condition is a \emph{boundary condition} at
the SUSY-breaking scale, not an RG fixed point.

\paragraph{Seed triples.}
For a triple $(0, \zch_a, \zch_b)$ with one massless member,
the energy-balance condition becomes:
$K = (\zch_a^2 + \zch_b^2)/(\zch_a + \zch_b)^2$.
Setting $r = \zch_a/\zch_b$ and solving $K = 2/3$:
$r^2 + 1 = 2(r + 1)^2/3$, giving $r = 2 - \sqrt{3}
\approx 0.268$.
This is a universal ratio: \emph{any} seed triple satisfying
$K = 2/3$ must have $\zch_{\rm small}/\zch_{\rm large}
= 2 - \sqrt{3}$.
\begin{itemize}
  \item \textbf{Meson seed} $(0, \pi, D_s)$:
    $\zch_\pi/\zch_{D_s} = 0.266$ ($0.6\%$ off),
    $K = 0.668$ ($0.18\%$).
  \item \textbf{Quark seed} $(0, s, c)$:
    $\zch_s/\zch_c = 0.271$ ($1.2\%$ off),
    $K = 0.664$ ($0.35\%$).
  \item \textbf{Leptons} $(e, \mu, \tau)$:
    $K = 0.666661$ ($0.001\%$)---an ``almost-seed''
    ($\zch_e = 0.72\MeV^{1/2}$ is small but nonzero).
\end{itemize}

\paragraph{The meson triple $(-\pi, D_s, B)$.}
With signed charges $\zch_\pi < 0$, $\zch_{D_s} > 0$,
$\zch_B > 0$:
\begin{equation}
  K(-\pi, D_s, B) = \frac{\zch_\pi^2 + \zch_{D_s}^2 + \zch_B^2}
  {(-\zch_\pi + \zch_{D_s} + \zch_B)^2}
  = 0.6674\,,
\end{equation}
$0.10\%$ from $2/3$---the second most precise triple known.
Among 660 meson combinations (11 pseudoscalars, 4 sign patterns),
only $(-\pi, D_s, B)$ and $(-\pi, D_s, B_s)$ lie within $0.2\%$.

The pseudoscalar triple $K(\pi, K, \eta) = 0.36$ is far from
$2/3$---and this is \emph{expected}, not anomalous.
The GOR relation~\cite{GellMannOakesRenner:1968}
$M_\pi^2 = B_0(m_u + m_d)$ gives boson mass \emph{squared}
linear in quark mass; composing with $m_q \propto \zch^2$ gives
$M_\pi^2 \propto (\zch_u^2 + \zch_d^2)$, a sum of two
squares---not a single $\zch^2$.
Heavy mesons have $m_{\rm meson} \approx m_Q + \Lambda_{\rm QCD}$
(HQET), so $m_{\rm meson} \approx m_Q \propto \zch_Q^2$: the
charge-squared law for the heavy quark propagates to the heavy
meson mass.
The pion enters with a \emph{negative} sign because it approaches
$m = 0$ from below in the chiral limit (a pseudo-Goldstone
boson).

\paragraph{Sumino protection.}
The charged-lepton $K = 0.666661$ ($0.001\%$) is too precise to
survive QED corrections without protection.
Sumino~\cite{Sumino:2009} showed that a gauged $\UU(3)_F$ with
$\alpha_{\rm em} = \alpha_F/4$ cancels the dangerous
$\alpha\,m_i\log m_i^2$ corrections, placing family gauge bosons
at $10^2$--$10^3\TeV$.


%======================================================================
\section{The emergent $\mathcal{N}{=}2$ structure}
\label{sec:N2}
%======================================================================

\subsection{What ``emergent $\mathcal{N}{=}2$'' means}

% Cross-reference to companion paper
The detailed analysis of the emergent $\mathcal{N}{=}2$ structure,
including the Seiberg--Witten curve at the Koide point and the
Cayley--Hamilton cubic, is developed in a companion
paper~\cite{companion:N2}; here we summarise the key features.

In genuine $\mathcal{N}{=}2$ theories~\cite{SeibergWitten:1994},
each matter field sits in a hypermultiplet: a pair of
$\mathcal{N}{=}1$ chirals $(\Phi, \tilde\Phi)$ related by
$\SU(2)_R$.
In our $\mathcal{N}{=}1$ theory, the confined spectrum
automatically produces two chiral multiplets per family:
\begin{equation}
  (M_i, \;\adj(M)_i)
  \qquad\text{per family } i = 1, 2, 3\,,
\end{equation}
with complementary quantum numbers:
$M^{1,2} \sim (\mathbf{2}, -1/2)$ while
$\adj(M)^{1,2} \sim (\mathbf{2}, +1/2)$.
This has the \emph{mathematical structure} of a
hypermultiplet---emergent from confinement plus the adjugate
identity, not from UV symmetry.


\subsection{Breaking the doublet: the quartic}

An auxiliary adjoint $\Phi$ (transforming as $\mathbf{8}$ of
$\SU(3)_F$) with mass $\mu_\Phi$ and Yukawa coupling
$W \supset \sqrt{2}\,g\,\tilde{Q}\,\Phi\,Q$~\cite{HMS:1997}
is integrated out via its $F$-term
$F_\Phi = \sqrt{2}\,g\,\tilde{Q}Q + \mu_\Phi\Phi = 0$,
giving $\Phi = -\sqrt{2}\,g\,\tilde{Q}Q/\mu_\Phi$.
Substituting back:
\begin{equation}
  W_{\rm quartic} = -\frac{g^2}{\mu_\Phi}
    \left[\Tr M^2
    - \frac{1}{3}(\Tr M)^2\right].
  \label{eq:quartic}
\end{equation}
The adjoint $\Phi$ is NOT a physical particle in the sBootstrap:
it is a UV field with mass $\mu_\Phi$ that is integrated out
below that scale, leaving only the quartic
coupling~\eqref{eq:quartic} as its low-energy footprint.
This is standard EFT decoupling---no different from the $W$
boson producing the Fermi four-fermion interaction below
$M_W$---and requires no special assumption beyond $\mu_\Phi$
being large compared to $\Lambda_F$.
As $\mu_\Phi \to \infty$, the quartic vanishes and pure
$\mathcal{N}{=}1$ SQCD is recovered.
At finite $\mu_\Phi$, the quartic couples the diagonal meson
entries to each other, mixing the direct and inverse branches.

\paragraph{The Yukawa coupling.}
In genuine $\mathcal{N}{=}2$
SUSY~\cite{Fayet:1976,SeibergWitten:1994},
the Yukawa coupling equals the gauge coupling:
$y = \sqrt{2}\,g$.
The UV theory here is $\mathcal{N}{=}1$, where $y$ is a priori
free.
We \emph{adopt} $y = \sqrt{2}\,g$ as the simplest choice
consistent with the emergent doublet structure---it is the value
that would hold if the $(M_i, \adj(M)_i)$ pairing were an exact
$\mathcal{N}{=}2$ hypermultiplet.
This is an \emph{assumption}, not a derivation from the
$\mathcal{N}{=}1$ dynamics.
If it holds even approximately, the Yukawa hierarchy is not a
hierarchy of couplings but of VEVs (via the Seiberg seesaw).


\subsection{Both Higgs doublets from one meson}

The structural consequence of the
$\mathcal{N}{=}2 \to \mathcal{N}{=}1$ breaking is that the
doublet $(M_i, \adj(M)_i)$ splits into components with
\emph{different} electroweak quantum numbers:
\begin{align}
  M^{1,2} &\sim (\mathbf{2}, -1/2) = H_d
    \qquad\text{(down-type Higgs)}\,,\\
  \adj(M)^{1,2} &\sim (\mathbf{2}, +1/2) = H_u
    \qquad\text{(up-type Higgs)}\,.
\end{align}
\textbf{Both MSSM Higgs doublets emerge from a single meson
matrix.}
This is not a postulate but a consequence of the adjugate
structure: $M$ and $\adj(M)$ automatically carry conjugate
quantum numbers.

There is no independent $\tan\beta$ parameter.
In the MSSM, $\tan\beta = \langle H_u\rangle/\langle H_d\rangle$
is free; here $H_d = M^{1,2}$ and $H_u = \adj(M)^{1,2}$ are both
determined by the same meson matrix and constrained by
$\det M = \Lambda^6$: their ratio is fixed by Seiberg dynamics,
not adjustable.
The ``superpartners'' are the mesons themselves, split from
mesinos by $f_\pi$; there is no second SUSY at the TeV scale.
For mixed up/down charge triples, $K(m_c, m_b, m_t) = 2/3$
involves both Higgs couplings symmetrically, consistent with
equal VEVs ($\adj(M) \cdot M = \det M$ on the constraint
surface).


\subsection{The critical scale and self-completion}

The energy-balance vacuum exists only above
\begin{equation}
  \mu_{\Phi,\rm crit} \approx 4.5\TeV\,.
  \label{eq:mucrit}
\end{equation}
Under 5-loop SMDR running~\cite{MartinRobertson:2019,PDG:2024}, $K^{-1}(d,s,b)$
passes through $2/3$ at $Q \approx 280\TeV$; the quartic must be
large enough to sustain the mixed vacuum, requiring
$\mu_\Phi \lesssim 4.5\TeV$.

The model self-completes: above $\Lambda_F$, the theory is
$\SU(3)_F \times \text{SM}$ with preons, asymptotically free
to the Planck scale.
The only additional scale is the family gauge boson mass at
$10^2$--$10^3\TeV$~\cite{Sumino:2009}.
There are no stops, sbottoms, gluinos, or conventional SUSY
partners: the ``SUSY'' refers to composite-sector structure,
not a TeV-scale extension.


%======================================================================
\section{SUSY breaking: seeds, bloom, and bions}
\label{sec:susy_breaking}
%======================================================================

\subsection{The O'Raifeartaigh model and the seed}

The O'Raifeartaigh superpotential~\cite{ORaifeartaigh:1975}
with three chiral superfields:
\begin{equation}
  W_{\rm OR} = f\,\Phi_0 + m\,\Phi_1\Phi_2
  + g\,\Phi_0\Phi_1^2\,,
  \label{eq:W_OR}
\end{equation}
breaks SUSY because the $F$-terms
\begin{align}
  F_0 &= f + g\,\Phi_1^2\,,\quad
  F_1 = m\,\Phi_2 + 2g\,\Phi_0\,\Phi_1\,,\quad
  F_2 = m\,\Phi_1
\end{align}
cannot simultaneously vanish ($F_0 = 0$ requires
$\Phi_1^2 = -f/g$, but then $F_2 = m\Phi_1 \neq 0$).
The fermionic mass matrix at
$\langle\Phi_0\rangle = v$,
$\langle\Phi_1\rangle = \langle\Phi_2\rangle = 0$:
\begin{equation}
  \mathcal{M}_F = \begin{pmatrix}
    0 & 0 & 0 \\ 0 & 2gv & m \\ 0 & m & 0
  \end{pmatrix}\,,
  \label{eq:fermion_mass_matrix}
\end{equation}
has eigenvalues $m_\pm = \sqrt{g^2v^2 + m^2} \pm gv$ and
$m_0 = 0$ (the goldstino).
At the seed condition $gv/m = \sqrt{3}$:
$m_- = (2-\sqrt{3})\,m$ and $m_+ = (2+\sqrt{3})\,m$.
Computing $K$:
\begin{equation}
  K(0, m_-, m_+) = \frac{m_- + m_+}{(\sqrt{m_-}+\sqrt{m_+})^2}
  = \frac{4m}{6m} = \frac{2}{3}\,,
\end{equation}
exactly.
K\"ahler stabilisation at the pseudo-modulus pole with
$\Lambda_K = m$ requires:
\begin{equation}
  c = -\frac{1}{12}\,.
  \label{eq:c_12}
\end{equation}

\paragraph{Why this is ``mild'' SUSY breaking.}
The fermion masses are $\mathcal{O}(m)$, while the SUSY-breaking
scale is $f$.
In the ISS dictionary~\cite{ISS:2006}, $g \to h$,
$m \to h\mu_{\rm ISS}$, $v \to X$, so the seed condition becomes
$\langle X\rangle/\mu_{\rm ISS} = \sqrt{3}$.
The ``mildness'' is the statement that $f \ll m^2$: in the
sBootstrap, $f = f_\pi^2 \approx (92\MeV)^2$ while
$m = \mu/4 \approx 471\MeV$, giving $f/m^2 \approx 0.038$.
This is why mesons and leptons appear at the same energy scale:
for the electron family
$m_{e,\rm scalar} = \sqrt{m_e^2 + f_\pi^2} \approx f_\pi
= 92\MeV$, while for the tau the splitting is negligible
($m_{\tau,\rm scalar} \approx m_\tau$).


\subsection{The bloom: from seed to physical spectrum}

The seed is $(0, m_-, m_+)$ with one vanishing mass.
The physical lepton spectrum $(m_e, m_\mu, m_\tau)$ has
$K = 0.666661$ but no vanishing mass.
The \textbf{bloom} is the nonperturbative process that transforms
the seed into the physical spectrum.

In the charge parametrisation
$\zch_k = z_0\bigl(1 + \sqrt{2}\cos(\delta + 2\pi k/3)\bigr)$,
the seed corresponds to $\delta = 3\pi/4$ (one charge vanishes).
The bloom is a rotation in $\delta$ at fixed scale.
But the seed is an \emph{unstable fixed point}: any
R-symmetry-breaking perturbation pushes $K$ away from $2/3$.
Shih~\cite{Shih:2008} showed that in O'Raifeartaigh models with
fields of R-charge $\neq 0, 2$, the one-loop Coleman--Weinberg
potential gives $m_X^2 < 0$, destabilising the pseudo-modulus
away from the origin---precisely the mechanism needed to push the
seed toward the physical spectrum.
The bloom must therefore be a nonperturbative rearrangement,
not a smooth deformation---it is driven by the instanton
potential described below.


\subsection{The instanton potential and $\mathbb{Z}_6$}

The anomalous $\UU(1)_A$ of $\SU(\Nc)$ SQCD with $\Nf$ flavours
has anomaly coefficient $2\Nf$~\cite{Seiberg:1994bz,Tong:SUSY}.
Instantons break $\UU(1)_A$ to $\mathbb{Z}_{2\Nf}$;
for $\Nf = 3$ the surviving discrete symmetry is
$\mathbb{Z}_6$.
(The anomaly-free $\UU(1)_R$ with $R[Q] = (\Nf - \Nc)/\Nf = 0$
is exact and continuous for $\Nc = \Nf$.)

In the charge parametrisation
$\zch_k = z_0(1 + \sqrt{2}\cos(\delta + 2\pi k/3))$,
every physical observable has the $\mathbb{Z}_3$ permutation
symmetry $\delta \to \delta - 2\pi/3$.
The effective potential is therefore a Fourier series in
$\cos(3n\delta)$; the leading instanton term is
\begin{equation}
  V_{\rm inst} \propto -\cos(3\delta)\,,
  \label{eq:cos3delta}
\end{equation}
with three minima at $\delta = 0, 2\pi/3, 4\pi/3$.
Higher harmonics $\cos(6\delta), \cos(9\delta), \ldots$ are
present but suppressed by $(\Lambda/\mu)^{b_0}$ per
instanton.
Note that $\delta$ parametrises mass \emph{ratios}
(the Cartan direction) and is independent of the $\UU(1)_A$
phase of $\det M$; the $\mathbb{Z}_3$ permutation symmetry
constraining $V(\delta)$ is a geometric property of the charge
lattice, while $\mathbb{Z}_6$ acts on the axial phase.
On the constraint surface $\det M = \Lambda^6$, holomorphic
invariants like $(\det M)^3$ are constant;
the $\delta$-dependent potential arises from K\"ahler corrections
and off-shell fluctuations around the constraint.
The physical lepton phase $\delta_{\rm phys}$ lies near one
of these minima.
At strong coupling, $V_0^{1/4} \sim 65$--$85\MeV \sim f_\pi$.


\subsection{The $\zch_b$ relation}

The following empirical relation holds to remarkable precision:
\begin{equation}
  \boxed{\zch_b = 3\zch_s + \zch_c\,,}
  \label{eq:v0_doubling_intro}
\end{equation}
with $3\zch_s + \zch_c = 64.67$ vs.\ $\zch_b = 64.68$
($0.01\%$).

A \emph{possible} dynamical origin comes from bion physics.
On $\mathbb{R}^3 \times S^1$, $\SU(\Nc)$ instantons
fractionalise into $\Nc$
monopole-instantons~\cite{Unsal:2008}.
Monopole--anti-monopole (bion) cross-terms contribute to the
K\"ahler potential:
\begin{equation}
  K_{\rm bion} = \frac{\zeta^2}{\Lambda^2}
  \Bigl|\sum_k s_k \zch_k\Bigr|^2\,,
  \label{eq:Kbion}
\end{equation}
where $s_k = \pm 1$ are topological signs.
Organising into seed ($S_{\rm seed} = \zch_s + \zch_c$) and
bloom ($S_{\rm bloom} = -\zch_s + \zch_c + \zch_b$) channels,
the second term is a perfect square vanishing when
$S_{\rm bloom} = 2\,S_{\rm seed}$, reproducing
eq.~\eqref{eq:v0_doubling_intro}.

\paragraph{Caveat.}
The \"Unsal bion analysis~\cite{Unsal:2008} was developed for
QCD(adj) (adjoint fermions on $\mathbb{R}^3 \times S^1$), not
for SQCD with fundamental matter.
Extending it to the present theory ($\Nc = \Nf = 3$ with
fundamentals) is plausible but unproven.
Equation~\eqref{eq:v0_doubling_intro} should therefore be
regarded as an \emph{empirical relation} ($0.01\%$) with a
suggestive but not yet rigorous bion motivation.
The predicted $m_b = \zch_b^2 = 4183\MeV$
(PDG: $4183 \pm 30\MeV$).

\paragraph{Bloom dynamics: frozen-sign vs.\ adiabatic.}
The adiabatic bion potential
$V_{\rm ad} = |\sum s_k \zch_k|^2$ is exactly flat on the
energy-balance manifold (an algebraic identity: $\sum s_k
|\zch_k| = 3v_0$ is constant).
It cannot select the phase $\delta$.
When topological signs are \emph{frozen} at the seed
configuration (where one amplitude vanishes), $V_{\rm frozen}
= (\sum |\zch_k|)^2$ provides a restoring force opposing the
bloom.
The equilibrium of $\cos(3\delta)$ (driving) and
$V_{\rm frozen}$ (restoring) determines the physical phase.

\paragraph{\"Unsal continuity caveat.}
Connecting the $\mathbb{R}^3 \times S^1$ analysis (where
monopole-instantons are under semiclassical control) to
$\mathbb{R}^4$ (where confinement is fully non-perturbative)
requires the \"Unsal continuity
conjecture~\cite{Unsal:2008}, supported by lattice evidence but
unproven for theories with matter.
The bion relation~\eqref{eq:v0_doubling_intro} should be regarded as
a prediction of the semiclassical regime whose validity at
$\Lambda \sim 300\MeV$ is assumed.

\paragraph{Numerical bloom analysis.}
The physical lepton phase $\delta_{\rm phys} = 132.73^\circ$
(from PDG masses) lies $12.73^\circ$ above the nearest instanton
minimum at $120^\circ$ and $2.27^\circ$ below the seed at
$135^\circ$.
The shift $\varepsilon = \delta_{\rm seed} - \delta_{\rm phys}
= 0.040\,\text{rad}$ generates the electron mass through
$m_e \approx \mu\varepsilon^2/6 = 0.49\MeV$ ($96\%$ of the
observed value; exact parametrisation gives $0.511\MeV$).

A one-loop Coleman--Weinberg computation of $V_{\rm CW}(\delta)$
on the constraint surface---including the quartic from the
adjoint mass, the off-diagonal meson loops, and the SUSY-breaking
$F$-terms---finds that \emph{all} perturbative contributions
minimise at $\delta = 120^\circ$ (the ``democratic'' point where
$\zch_1 = \zch_2$), not at $\delta_{\rm phys}$.
The physical deviation $\alpha_{\rm lep} = \delta_{\rm phys}
- 120^\circ = 12.73^\circ$ is a \emph{flat direction} of the
perturbative effective potential: the tree-level quartic gives
$V_{\rm quartic} \propto \mathrm{Var}(1/\zch_i^2)$ (minimised
at the democratic point), and the CW correction has
$\mathrm{STr}\,\mathcal{M}^4 \propto V_{\rm quartic}$ (same
angular profile, no shift).
The angle $\alpha_{\rm lep}$ therefore requires non-perturbative
K\"ahler dynamics, confirming that the Cartan angle and the bloom dynamics
are facets of the scalar K\"ahler metric problem.


\subsection{$F$-terms are lepton masses}

When electroweak symmetry breaks, the $F$-term system of the
confined theory becomes overconstrained: 18~real equations for
12~variables.
SUSY necessarily breaks, with:
\begin{equation}
  |F_{2j}| = m_j\,,\qquad j = e, \mu, \tau\,.
  \label{eq:Fterms_leptons}
\end{equation}
The goldstino is predominantly the $\tau$-mesino (largest
$F$-term).
The gravitino mass
$m_{3/2} = |F|/(\sqrt{3}\,M_{\rm Pl})
\approx 4 \times 10^{-16}\GeV$
is negligibly small.

\paragraph{The scalar--fermion splitting.}
The universal Volkov--Akulov splitting is
$m^2_{\rm scalar} = m^2_{\rm fermion} + f_\pi^2$.
For the pion--muon pair:
$m_\pi^2 = m_\mu^2 + f_\pi^2 = 11164 + 8501 = 19665\MeV^2$,
giving $m_\pi = 140.2\MeV$ (PDG: $139.6\MeV$, $0.5\%$).
% REVISED: REF-007
This is the foundational Volkov--Akulov check: the pion (scalar
meson) and the muon (fermionic mesino) are partners in a chiral
multiplet, split by $f_\pi$.


\subsection{Connection to the Seiberg dynamics}

The O'Raifeartaigh model is not an independent postulate but the
SUSY-breaking \emph{deformation} of the Seiberg vacuum.
The three constraining mechanisms work in concert:
\begin{enumerate}
  \item \textbf{Seiberg constraint} $\det M = \Lambda^6$
    (Sec.~\ref{sec:seiberg}): fixes the product
    $a_i \cdot d_i = \mu\nu$ (family-independent).
  \item \textbf{O'Raifeartaigh seed} $gv/m = \sqrt{3}$
    (this section): fixes the charge ratio within a seed triple
    to the algebraic value $r = 2 - \sqrt{3}$.
  \item \textbf{Instanton + bion bloom}: the $\cos(3\delta)$
    potential redistributes charge from the seed to the physical
    spectrum, while the bion restoring force constrains one mass
    combination ($\zch_b = 3\zch_s + \zch_c$).
\end{enumerate}
The quartic from the adjoint mass (Sec.~\ref{sec:N2}) couples
the direct and inverse branches, but whether this coupling is
sufficient to produce the full waterfall structure remains
open (Sec.~\ref{sec:inverse}).
The overall picture is a hierarchy of constraints: the Seiberg
constraint gives a one-dimensional family of vacua (parametrised
by $X$), the O'Raifeartaigh seed selects a point on this family
(the $\sqrt{3}$ condition), and the bloom dynamics perturbs it
to the physical spectrum.


%======================================================================
\section{The inverse sector and chain}
\label{sec:inverse}
%======================================================================

\subsection{The dual energy-balance on $(d, s, b)$}

Evaluating on inverse masses
($\zch_j \to 1/\zch_j$):
\begin{equation}
  K^{-1}(d,s,b) = \frac{1/\zch_d^2 + 1/\zch_s^2 + 1/\zch_b^2}
  {(1/\zch_d + 1/\zch_s + 1/\zch_b)^2}
  = 0.6652\,,
\end{equation}
$0.22\%$ from $2/3$, using PDG 2024 $\overline{\rm MS}$
masses~\cite{PDG:2024}.
Under RG running, both triples evolve.
Using 5-loop SMDR running with the Martin--Robertson
code~\cite{PDG:2024}, we scan from $1\GeV$ to $M_{\rm Planck}$
with 391 points, computing $\overline{\rm MS}$ quark masses at
each scale with the frozen-mass prescription (virtual particles
below their mass threshold contribute their pole mass, not a
running mass).
Error propagation uses PDG fractional uncertainties, numerical
$\partial K/\partial m_i$, added in quadrature.

The inverse down-type triple improves monotonically toward
the UV:
\begin{center}
\begin{tabular}{lll}
\hline
$Q$ & $K^{-1}(d,s,b)$ & deviation \\
\hline
$M_Z$ & $0.66750$ & $+0.126\%$ \\
$1\TeV$ & $0.66719$ & $+0.078\%$ \\
$10\TeV$ & $0.66695$ & $+0.042\%$ \\
$100\TeV$ & $0.66675$ & $+0.012\%$ \\
$\mathbf{280\TeV}$ & $\mathbf{0.66667}$ & $\mathbf{0.000\%}$ \\
$10^6\GeV$ & $0.66657$ & $-0.014\%$ \\
$M_{\rm Planck}$ & $0.66539$ & $-0.191\%$ \\
\hline
\end{tabular}
\end{center}
The crossing $K^{-1}(d,s,b) = 2/3$ occurs at
$Q = 280\TeV$ (5-loop SMDR, interpolated).
A second inverse triple $K^{-1}(s,c,-t)$ crosses $2/3$
at $Q = 82\TeV$.
Both share the strange quark as pivot, forming an
\textbf{inverse chain}: $(d,s,b)$--$(s,c,-t)$.

Key $\sigma$ values at the $\overline{\rm MS}$ scale:
$K(e,\mu,\tau)$ has $\sigma = 0.001\%$ (153$\sigma$ from exact
$2/3$);
$K^{-1}(s,c,t)$ has $\sigma = 0.13\%$ ($1.0\sigma$ from $2/3$);
$K^{-1}(d,s,b)$ has $\sigma = 0.26\%$ ($0.5\sigma$ from $2/3$).
The lepton Koide is measured; the quark inverse Koide values are
consistent with $2/3$ but with larger uncertainties from the
quark mass determinations.

An exhaustive scan of all 672 configurations of six quarks
% REVISED: REF-008
($\binom{6}{3} = 20$ mass triples $\times$ $2^3 - 1 = 7$
non-trivial sign patterns $\times$ direct and inverse $= 280$;
plus $\binom{5}{3} \times 7 \times 2 = 140$ five-quark subsets;
plus mixed-sector triples: total 672)
finds 31 sub-$0.5\%$ tuples and 27 zero crossings.
Cross-sector tuples (e.g.\ $u, \mu, \tau$) are transient
zero-crossings under RG running and are discarded as
unphysical.


\subsection{The interleaving theorem}

Given a single set of three charges $\zch_1 < \zch_2 < \zch_3$,
the Seiberg dynamics produces the \textbf{direct branch}
$a_i = \mu\zch_i^2$ and the \textbf{inverse branch}
$d_i = \nu/\zch_i^2$, with opposite orderings.

\begin{theorem}[Interleaving]
\label{thm:interleave}
The six masses form the alternating ladder
\begin{equation}
  d_1 > a_3 > d_2 > a_2 > d_3 > a_1
\end{equation}
if and only if
$\max(\zch_1^2\zch_3^2, \zch_2^4) < \nu/\mu
< \zch_2^2\zch_3^2$.
\end{theorem}

\begin{proof}
Each inequality constrains $\nu/\mu$.
$d_1 > a_3$ requires $\nu/\zch_1^2 > \mu\zch_3^2$,
i.e.\ $\nu/\mu > \zch_1^2\zch_3^2$.
$a_3 > d_2$ requires $\mu\zch_3^2 > \nu/\zch_2^2$,
i.e.\ $\nu/\mu < \zch_2^2\zch_3^2$.
The remaining four inequalities are weaker or equivalent.
The tightest lower bound is
$\max(\zch_1^2\zch_3^2, \zch_2^4)$; the upper is
$\zch_2^2\zch_3^2$.
\end{proof}

The product $a_i \cdot d_i = \mu\nu$ is family-independent
(from $\det M = \Lambda^6$)---the geometric mean constraint:
\begin{equation}
  a_i \cdot d_i = \mu\nu = \text{const}\,,
  \qquad\forall\; i\,.
  \label{eq:productrule}
\end{equation}

\paragraph{Quantitative status.}
The interleaving theorem is \emph{proven} as abstract mathematics
on any set of three charges.
However, the identification of the six physical quark masses with
the six rungs faces a quantitative obstruction.
If we take the down-type charges $(\zch_d, \zch_s, \zch_b)$ as
the direct branch and compute inverse masses from
$d_i = \nu/\zch_i^2$ with $\nu = \zch_d^2 \cdot m_t$ (requiring
$d_1 = m_t$), then $d_2 = \nu/\zch_s^2 = 8638\MeV$ versus
$m_c = 1273\MeV$---a factor $6.8$ discrepancy.
This is the same factor appearing in the holomorphic $m_c/m_s$
ratio (Sec.~\ref{sec:composites}): the K\"ahler $Z$-factors
$Z_c/Z_s \approx 1.37$ ($1.37^2 \approx 1.9$) reduce but do
not eliminate the gap.
\textbf{The waterfall cannot reproduce the up-type spectrum from
down-type charges alone without K\"ahler corrections.}
This is tied to the master obstruction
(Problem~1 in Sec.~\ref{sec:open}).


\subsection{The chain structure}

The quark mass spectrum exhibits overlapping triples:
\begin{equation}
  \underbrace{(t,b,c)}_{\text{Fermi}} \;-\;
  \underbrace{(b,c,-s)}_{\text{transition}} \;-\;
  \underbrace{(c,s,0)}_{\text{seed}} \;-\;
  \underbrace{(s,0,m_u{+}m_d)}_{\text{lightest}}\,.
  \label{eq:chain}
\end{equation}

\begin{center}
\begin{tabular}{llll}
\hline
Link & Triple & $K$ & Status \\
\hline
1 & $(+t, +b, +c)$ & $0.6693$ ($0.40\%$) & Fermi side \\
2 & $(+b, +c, -s)$ & $0.6747$ ($1.21\%$) & sign flip on $s$ \\
3 & $(+c, +s, 0)$ & $0.6646$ ($0.31\%$) & seed triple \\
4 & $(+s, 0, \zch_{u+d})$ & $0.6649$ ($0.28\%$) & lightest seed \\
\hline
\end{tabular}
\end{center}

\paragraph{Connection to the interleaving.}
The chain and the waterfall are logically \textbf{independent}
observations.
The interleaving theorem predicts an alternating pattern of six
masses; the chain is a sequence of overlapping triples.
Whether the chain \emph{arises from} the waterfall has not been
proven.
The quantitative discrepancy (factor $6.8$ in the up-type
spectrum) suggests that additional physics beyond the pure
interleaving is needed.

\paragraph{Alternative origins for the chain.}
Several mechanisms could produce the overlapping-triple pattern
independently of the waterfall:
\begin{enumerate}
  \item \textbf{SUSY breaking:} the O'Raifeartaigh seed
    generates $(0, m_-, m_+)$ with $K = 2/3$.
    Successive blooms at different scales could produce a
    sequence of seed triples whose overlap creates the chain.
  \item \textbf{Charge additivity:} the bion relation
    $\zch_b = 3\zch_s + \zch_c$ shows that $\zch_b$
    \emph{decomposes} into $\zch_c$ plus three $\zch_s$.
    If charges are additive (as expected for $\UU(1)_F$),
    overlapping decompositions naturally produce triples
    sharing members.
  \item \textbf{Kinematic constraint:} with six quarks ordered
    by mass, there are exactly four windows of three consecutive
    quarks.  The chain is simply \emph{all} windows; what
    requires explanation is not the chain itself but why
    $K \approx 2/3$ in each window.
  \item \textbf{Multiple $\UU(1)$ charges:} the Cartan of
    $\SU(3)_F$ has \emph{two} $\UU(1)$ generators; projecting
    onto different linear combinations could produce sign flips
    and different effective charges.
\end{enumerate}
At present we cannot distinguish between these origins.
The chain is an \emph{observed pattern} with quantitative support;
its dynamical explanation remains open.

\paragraph{The repeated seed map.}
The chiral tail is two iterations of the seed map
$\zch_{k-1} = r\,\zch_k$ with $r = 2 - \sqrt{3}$:
$\zch_s/\zch_c = 0.271$ and $\zch_{u+d}/\zch_s = 0.271$,
both within $1.1\%$ of $r \approx 0.268$.

\paragraph{Threshold decoupling and the chiral tail.}
The ``chiral'' appearance of the last two links is not accidental.
Below the charm threshold, heavy flavours have been integrated out
in the sequence $6 \to 5 \to 4 \to 3$ active quarks
(Table~\ref{tab:phases}).
In this regime the natural variables are those of chiral
perturbation theory: the GOR combination
$m_u + m_d$ (proportional to $M_\pi^2$) replaces the individual
up and down masses, and the vanishing charge $\zch = 0$ represents
the chiral limit.
The chain's tail is an EFT reorganisation of variables after
threshold decoupling, not a statement about six simultaneous
current masses at a common scale.
This also gives a structural reading of the coefficient 3 in the
bion relation $\zch_b = \zch_c + 3\zch_s$: three threshold
crossings ($t \to b \to c \to s$) separate the $b$-quark from
the strange-pivoted chiral sector.

\paragraph{Top completion.}
The top link can be \emph{derived} rather than fitted.
Solving $K[\zch_t, \zch_b, \zch_c] = 2/3$ for positive $\zch_t$
is a quadratic equation whose positive root gives the exact
completion formula:
\begin{equation}
  \zch_t^{\rm pred} = 2(\zch_b + \zch_c)
  + \sqrt{3}\,\sqrt{\zch_b^2 + 4\,\zch_b\zch_c + \zch_c^2}\,.
  \label{eq:top_completion}
\end{equation}
This is the unique positive real root; the negative root would
give a ``negative mass'' solution which is unphysical.
Numerically: $\zch_t^{\rm pred} = 410.5\MeV^{1/2}$, predicting
$m_t^{\rm pred} = \zch_t^2 = 168.5\GeV$, compared with
$m_t^{\rm pole} = 172.4\GeV$ ($2.3\%$ discrepancy).
This is a genuine \emph{prediction}: the top mass is determined
by the bottom and charm masses through the energy-balance
condition, with no free parameters.

\paragraph{Chain compression.}
The chain is thus \emph{compressed by three relations}:
\begin{align}
  \zch_s &\approx (2-\sqrt{3})\,\zch_c
    &&\text{(seed map)}\,,\nonumber\\
  \zch_{u+d} &\approx (2-\sqrt{3})\,\zch_s
    &&\text{(iterated seed)}\,,\nonumber\\
  \zch_b &\approx \zch_c + 3\zch_s
    &&\text{(bion / threshold)}\,,
  \label{eq:chain_compression}
\end{align}
plus the positive completion~\eqref{eq:top_completion} for
$\zch_t$.
Given $\zch_c$ as the single input, the remaining five charges
are determined to within $\sim\!2\%$.
The seed ratio $r = 2 - \sqrt{3}$ is an algebraic constant
emerging from the energy-balance condition for a zero-mass triple;
the bion coefficient $3$ counts threshold crossings; and the
completion formula is the unique positive root of a quadratic.
The entire six-quark mass spectrum is thus parametrised by a
single number ($\zch_c = 35.68\MeV^{1/2}$) plus four algebraic
relations.

\paragraph{The cross-sector product rule and the Cabibbo angle.}
If leptons and down quarks share the same three family charges
on opposite branches, the product rule $a_i \cdot d_i = \mu\nu$
(eq.~\ref{eq:productrule}) gives a \emph{cross-sector}
prediction at the holomorphic level:
\begin{equation}
  m_e \cdot m_b^{\rm hol} = m_\mu \cdot m_s^{\rm hol}
  = m_\tau \cdot m_d^{\rm hol}\,.
  \label{eq:cross_sector}
\end{equation}
This implies $m_d^{\rm hol}/m_s^{\rm hol} = m_\mu/m_\tau$,
and hence a holomorphic prediction for the Cabibbo angle via the
Gatto--Sartori--Tonin relation~\cite{GattoSartoriTonin:1968}:
\begin{equation}
  |V_{us}|_{\rm hol} = \sqrt{m_\mu/m_\tau} = 0.244\,,
  \label{eq:Vus_hol}
\end{equation}
compared with the observed $|V_{us}| = 0.224$ ($8.7\%$ off).
The K\"ahler correction needed is $Z_d/Z_s = 1.18$.

However, the full product rule~\eqref{eq:cross_sector} is
\textbf{badly violated}: $m_e \cdot m_b = 2138\MeV^2$ vs.\
$m_\mu \cdot m_s = 9869\MeV^2$ (factor~4.6).
This means the simple ``same charges, opposite branches''
identification requires large, generation-dependent K\"ahler
corrections across the lepton--quark boundary---far beyond the
mild ($< 1\%$) corrections needed within the lepton sector alone.
The cross-sector product rule is therefore an \emph{order of
magnitude} test of the framework, not a precision prediction;
its failure at the holomorphic level is the most concrete
manifestation of the K\"ahler metric obstruction.

\subsection{The CKM matrix: an open problem}

The holomorphic cross-sector prediction
$|V_{us}| = \sqrt{m_\mu/m_\tau} = 0.244$ is $8.7\%$ from the
observed value---not a precision result, but in the right
ballpark, suggesting the charge identification is qualitatively
correct with K\"ahler corrections of order $Z_d/Z_s \approx 1.2$.

The remaining CKM elements ($|V_{cb}|$, $|V_{ub}|$,
$\delta_{\rm CP}$) are \textbf{not predicted}.
Empirical patterns have been noted: the charge ratio
$R = \zch_b^2/\zch_c^2 = (5+\sqrt{21})/2$ satisfies the
self-dual equation $R^2 - 5R + 1 = 0$, and
$R^3 = 55 + 12\sqrt{21}$ is the fundamental unit of the real
quadratic field $\mathbb{Q}(\sqrt{21})$.
But these are \emph{observations about the data}, not derivations
from the theory.
Until the off-diagonal meson couplings and the K\"ahler metric
are computed, the CKM mixing matrix remains the most significant
gap in the framework.

The CP phase satisfies a suggestive $N = 3$ consistency condition:
$\cos^2(2\delta) = N/(N+4) = N/(2N+1)$ holds only for $N = 3$
(since $N + 4 = 2N + 1$ gives $N = 3$).
Whether this is structural or coincidental is not resolved.


%======================================================================
\section{Generation uniqueness and $\SO(32)$}
\label{sec:SO32}
%======================================================================

\subsection{Why five quarks in the mesons}

The SM has six quarks but the s-confining $\SU(3)_F$ requires
$\Nf = \Nc = 3$.
The top quark ($m_t = 173\GeV \gg \Lambda_F$) is integrated out:
it is too heavy to participate in the confinement.
The remaining five quarks ($u,d,s,c,b$) decouple sequentially
as the RG scale decreases (Table~\ref{tab:phases}).
The Diophantine system counts light quarks of each type
(up-type and down-type) that participate in confinement,
selecting $(r, s) = (3, 2)$: three down-type and two up-type
light quarks, with $N = rs/2 = 3$ generations.

\subsection{The Diophantine system}

The sBootstrap counting~\cite{Rivero:2024epjc} imposes three
simultaneous conditions on the numbers of up-type ($s$) and
down-type ($r$) light quarks participating in confinement:
\begin{align}
  rs &= 2N && \text{(total flavour pairs)}\,,\nonumber\\
  \frac{r(r+1)}{2} &= 2N && \text{(symmetric composites)}\,,
  \nonumber\\
  r^2 + s^2 - 1 &= 4N && \text{(adjoint content)}\,.
\end{align}
The first equation counts the number of meson flavour pairs.
The second requires the symmetric representation $\mathbf{15}$ of
$\SU(5)_f$ to appear with the correct multiplicity.
The third is a consistency condition on the total DOF.
Together they have the unique solution $(N, r, s) = (3, 3, 2)$:
three generations, with three down-type and two up-type light
quarks.
The composites fill $\mathbf{24} \oplus \mathbf{15} \oplus
\overline{\mathbf{15}}$ of $\SU(5)_f$.

\subsection{Additional $N = 3$ uniqueness criteria}

% REVISED: REF-005 — criteria ranked by strength, independence clarified
Beyond the Diophantine system, four additional criteria each
select $N = 3$.
We rank them by the strength of the physical argument.
The five criteria arise from different mathematical structures
(number theory, algebraic geometry, representation theory,
CP-violation counting, and Casimir invariants).
They are algebraically independent in the sense that no one
implies another, though they all constrain the same physical
quantity $N$ and share the common feature of being rooted in
$\SU(N)$ representation theory.

\begin{enumerate}
  \item \textbf{Discriminant matching} [\emph{suggestive}]:
    the discriminant of the characteristic polynomial of $\SU(N)$
    gives $D_{\rm alg} = (N+2)^2 - 4$, while the dimension of
    the s-confining constraint surface is $D_{\rm dyn} = N(2N+1)
    = \dim(\Sp(2N))$.
    Requiring $D_{\rm alg} = D_{\rm dyn}$ gives
    $(N+2)^2 - 4 = N(2N+1)$, which has the unique solution
    $N = 3$ ($21 = 21$).
    The physical principle connecting the algebraic discriminant
    to the constraint-surface dimension is not established; this
    criterion is a numerical coincidence that may indicate a
    deeper structure.

  \item \textbf{Binomial identity} [\emph{suggestive}]:
    $\binom{3N}{N} = (N+1) \cdot N(2N+1)$ holds only at $N = 3$:
    $84 = 4 \times 21$.
    The number $84$ counts the $C_3$ polarisations of 11d SUGRA
    under $\SO(9)$ ($\Lambda^3(\mathbb{R}^9) = \mathbf{84}$),
    and 21 is $\dim(\SU(3)_F \text{ composites at } \Nf = \Nc =
    3)$.
    While suggestive of a connection between 11d geometry and
    4d confinement, this remains an arithmetic coincidence
    without a derivation from dynamics.

  \item \textbf{CP-phase consistency} [\emph{conditional}]:
    requiring $\cos^2(2\delta) = N/(N+4)$ and
    $\cos^2(2\delta) = N/(2N+1)$ simultaneously gives
    $N + 4 = 2N + 1$, which requires $N = 3$.
    This uses the observed CP phase $\delta_{\rm CP}$, which is
    not currently predicted by the framework; it is a
    consistency check rather than a selection criterion.

  \item \textbf{Casimir--Koide identity} [\emph{structural}]:
    for $\UU(N)$ with canonical normalisation, the energy-balance
    ratio $K_N = 2/N$ (Theorem~\ref{thm:n2} generalised).
    The mean-square weight $\langle J_3^2\rangle
    = N(N^2-1)/(12N) = (N^2-1)/12$ of the principal $\mathfrak{sl}(2)$
    embedding in $\mathfrak{sl}(N)$~\cite{Kostant:1959}
    (where $J_3$ is the weight operator in the fundamental
    representation).
    The condition $K_N = \langle J_3^2\rangle$ gives
    $N(N^2-1) = 24$, whose unique positive integer solution is
    $N = 3$.
    Equivalently, at $N = 3$ and only $N = 3$, the representation
    dimension $2j + 1 = N$ coincides with the algebra dimension
    $\dim\,\mathfrak{sl}(2) = 3$, giving
    $C_2/N = C_2/3 = \langle J_3^2\rangle = K = 2/3$.
\end{enumerate}

\subsection{$\SO(32)$ selection}

The ten-dimensional heterotic string has two anomaly-free gauge
groups: $\SO(32)$ and $E_8 \times E_8$.
The sBootstrap composite content
$\mathbf{24} \oplus \mathbf{15} \oplus \overline{\mathbf{15}}$
of $\SU(5)_f$ selects between them.

For $\SO(32)$:
$\SO(32) \supset \SU(16) \times \UU(1) \supset \SU(5) \times
\SU(3) \times \UU(1)^2$.
The adjoint $\mathbf{496}$ contains $\mathbf{24}$ and
$\mathbf{15} + \overline{\mathbf{15}}$---precisely the
sBootstrap content.
The $Q = 4/3$ states in the $\mathbf{24}$ are projected out by
the $\SU(3)_L$ decomposition of
Sec.~\ref{sec:composites}, eq.~\eqref{eq:meson_decomp}: this is
the cleaning that makes the $\mathbf{15}$ SM-compatible.

For $E_8 \times E_8$: every decomposition of $E_8$ under
$\SU(5)$ produces only the antisymmetric $\mathbf{10}$, never
the symmetric $\mathbf{15}$~\cite{Rivero:2024epjc}.
The sBootstrap therefore selects $\SO(32)$ over
$E_8 \times E_8$.

\begin{table}
\centering\small
\caption{$\SU(5)_f$ composite content and sBootstrap scorecard.}
\label{tab:SU5_scorecard}
\begin{tabular}{llll}
\hline
$\SU(5)_f$ rep & SM decomposition & Role & Status \\
\hline
$\mathbf{24}$ & $(\mathbf{8},\mathbf{1})_0 + \ldots$ & gauge + Higgs & \\
  & $+ (\mathbf{3},\mathbf{2})_{-5/6} + \text{c.c.}$ & $Q{=}4/3$ proj.\ & \\
$\mathbf{15}$ & $(\mathbf{1},\mathbf{3})_{-1}$ & sleptons & $\checkmark$ \\
  & $(\mathbf{3},\mathbf{2})_{1/6}$ & leptoquarks & $\approx$ \\
  & $(\mathbf{6},\mathbf{1})_{-2/3}$ & diquarks & $\times$ \\
\hline
\multicolumn{4}{l}{84 DOF: exact 16 (19\%), approx.\ 44 (52\%), obstr.\ 24 (29\%).} \\
\hline
\end{tabular}
\end{table}

In 11d SUGRA, $\mathbf{6}_C$ belongs to
$\mathrm{Sym}^2(9) = \mathbf{45} \subset g_{MN}$ (gravity
sector), not $\Lambda^3(9) = \mathbf{84} \subset C_3$ (gauge
sector), explaining the colour-sextet obstruction.


%======================================================================
\section{The effective Lagrangian}
\label{sec:lagrangian}
%======================================================================

\subsection{Scenario (a): sequential integration}

\paragraph{The UV theory.}
The starting point is $\SU(3)_F$ SQCD with $\Nf = 5$ massive
flavours $(u,d,s,c,b)$.
As the RG scale decreases below $m_b$ and $m_c$, heavy flavours
are integrated out sequentially (Table~\ref{tab:phases}).
At each threshold the dynamical scale matches:
$\Lambda_3^6 = \Lambda_4^5 m_c$,
$\Lambda_4^5 = \Lambda_5^4 m_b$.

\paragraph{The $\Nf = 4$ ISS window.}
Between $m_c$ and $m_b$, the theory has $\Nf = 4$ and sits at the
lower boundary of the ISS window~\cite{ISS:2006}
($\Nc {+} 1 \leq \Nf < \frac{3}{2}\Nc$, i.e.\ $4 \leq \Nf < 4.5$).
The IR-free magnetic dual has gauge group
$\SU(\Nf - \Nc) = \SU(1) = \text{trivial}$ (for $\Nf = 4$,
$\Nc = 3$); $\Nf = 4$ is a boundary case where the
magnetic gauge dynamics is absent but the rank-condition
SUSY breaking persists~\cite{ISS:2006}.
% REVISED: REF-010
The macroscopic superpotential is:
\begin{equation}
  W_{\rm ISS} = h\,\Tr\,\varphi\,\Phi\,\tilde\varphi
  - h\mu^2\,\Tr\Phi\,,
  \label{eq:W_ISS}
\end{equation}
where $\Phi_{ij}$ is the meson of the electric theory (now
classical in the magnetic description), $\varphi$ and
$\tilde\varphi$ are magnetic quarks, $h$ is the magnetic Yukawa,
and $\mu^2 \sim m\Lambda$ is the SUSY-breaking parameter.
The ISS variables map to sBootstrap composites:
$\Phi_{ij} \to M^i{}_j/\Lambda$ (normalised meson),
$h\mu^2 \to m_i\Lambda$ (preon mass $\times$ scale).

SUSY is broken by the rank condition: $F_\Phi \sim
\tilde\varphi\varphi - \mu^2\,\mathbf{1}$ has rank
$\Nf - \Nc = 1$, but $\mu^2\,\mathbf{1}$ has rank $\Nf = 4$,
so not all $F$-terms can vanish.
The metastability parameter
$\epsilon = \mu/\Lambda_m \sim \sqrt{m/\Lambda} \ll 1$ ensures
parametrically long lifetime.
At $\Nf = 3$ the magnetic gauge group is trivial ($\SU(0)$),
so the magnetic quarks disappear and eq.~\eqref{eq:W_ISS}
reduces to the s-confined superpotential~\eqref{eq:W_conf}---the
ISS window serves only as a \emph{transient} SUSY-breaking phase
during the RG descent from $\Nf = 4$ to $\Nf = 3$.

\paragraph{The $\Nf = 3$ confined theory.}
Below $m_c$, the theory s-confines with the effective
superpotential:
\begin{equation}
  W_{\rm conf} = X\bigl(\det M^{(3)} - B\bar{B} - \Lambda_3^6\bigr)
    + \Tr(\hat{m}\,M^{(3)})\,,
  \label{eq:W_conf}
\end{equation}
where $M^{(3)}$ is the $3 \times 3$ light-flavour meson matrix,
$X$ is the Lagrange multiplier enforcing the quantum constraint,
and $\hat{m} = \diag(m_u, m_d, m_s)$ are the light preon masses.

\paragraph{The full superpotential.}
Combining all sectors:
\begin{align}
  W_{\rm (a)} &= X\bigl(\det M^{(3)} - B\bar{B} - \Lambda_3^6\bigr)
    + \Tr(\hat{m}\,M^{(3)})
    \nonumber\\
  &\quad + c_3\,\frac{(\det M^{(3)})^3}{\Lambda_3^{18}}
    \quad\text{(three-instanton, \S\ref{sec:susy_breaking})}
    \nonumber\\
  &\quad + y_c\,H_u\,M^c{}_c + y_b\,H_d\,M^b{}_b
    \quad\text{(heavy-flavour Yukawas)}
    \nonumber\\
  &\quad + y_t\,H_u\,Q^t\bar{Q}_t
    \quad\text{(top: partial compositeness)}\,.
  \label{eq:W_scenario_a}
\end{align}
The first line is the s-confined sector.
The second is the $\cos(3\delta)$ instanton potential
(Sec.~\ref{sec:susy_breaking}).
The third couples the charm and bottom mesons to the Higgs
doublets $H_u = \adj(M)^{1,2}$ and $H_d = M^{1,2}$.
The fourth is the top quark Yukawa: the top is too heavy to
confine ($m_t \gg \Lambda_F$) and enters through partial
compositeness~\cite{Kaplan:1991dc}, coupling directly to the
elementary preon fields.
Note that the charm and bottom enter differently from the top:
they participate in the confinement at the $\Nf = 4, 5$ levels
but are integrated out before s-confinement at $\Nf = 3$, so
they contribute through the threshold matching
$\Lambda_3^6 = \Lambda_4^5 m_c$ rather than through the
confined superpotential directly.

\paragraph{The K\"ahler potential.}
\begin{equation}
  K_{\rm (a)} = \Tr(M^\dagger M) + |X|^2 + |B|^2 + |\bar{B}|^2
  - \frac{|X|^4}{12\mu^2}
  + \frac{\zeta^2}{\Lambda^2}
    \Bigl|\sum_k s_k \zch_k\Bigr|^2\,.
  \label{eq:Kahler_a}
\end{equation}
The first four terms are canonical.
The $-|X|^4/(12\mu^2)$ term is the negative quartic K\"ahler
correction that stabilises the pseudo-modulus at $gv/m = \sqrt{3}$
(Sec.~\ref{sec:susy_breaking}, eq.~\eqref{eq:c_12}).
The last term is the bion contribution to the K\"ahler
potential~\cite{Unsal:2008}, with topological signs
$s_k = \pm 1$.
The soft potential $V_{\rm soft} = f_\pi^2\,\Tr(M^\dagger M)$
gives universal splitting
$m^2_{\rm scalar} - m^2_{\rm fermion} = f_\pi^2$.

\paragraph{The scalar--fermion splitting.}
The soft SUSY-breaking potential is
$V_{\rm soft} = f_\pi^2\,\Tr(M^\dagger M)$,
giving a universal splitting
$m^2_{\rm scalar} - m^2_{\rm fermion} = f_\pi^2$ for every
chiral multiplet.
The supertrace
$\mathrm{STr}[\mathcal{M}^2] = 18\,f_\pi^2$
(from 9 meson multiplets $\times$ 2 real scalars each).

% REVISED: REF-002 — transparent parameter counting
\paragraph{Parameter counting.}
Table~\ref{tab:params} lists all inputs to the framework.
In the holomorphic sector, the lepton mass
\emph{ratios} are set by $\delta_{\rm lep}$ (the Cartan angle);
quark chain ratios by $\delta_{\rm quark}$
(a different Cartan angle, giving the seed ratio
$\zch_s/\zch_c \approx 2 - \sqrt{3}$).
The absolute scale $\mu$ converts charge ratios to masses.
$\Lambda_F$ is fixed by matching to $\alpha_s(M_Z)$;
$\mu_\Phi$ is empirical ($\lesssim 4.5\TeV$); and the $Z_i$
are the master unknowns (Problem~1 in Sec.~\ref{sec:open}).
The mapping to physical masses also requires $Z_i$,
$\mu_\Phi$, and the bloom dynamics---none of which are derived
from $\delta$ and $\mu$ alone.

\begin{table}
\centering\small
\caption{Framework inputs and their roles.}
\label{tab:params}
\begin{tabular}{lll}
\hline
Input & Sets & Status \\
\hline
$\delta$ (Cartan) & mass ratios & fitted ($2/9$) \\
$\mu$ (scale) & abs.\ masses & fitted ($1883\MeV$) \\
$\Lambda_F$ & composites & $\alpha_s(M_Z)$ \\
$y{=}\sqrt{2}g$ & Yukawa--gauge & assumed ($\mathcal{N}{=}2$) \\
$\mu_\Phi$ & $\mathcal{N}{=}2{\to}1$ & ${\lesssim}\,4.5\TeV$ \\
$m{\propto}\zch^2$ & mass law & postulate \\
$Z_i$ (K\"ahler) & hol.\ $\to$ phys.\ & \textbf{unknown} \\
\hline
\multicolumn{3}{p{\columnwidth}}{%
\textit{Continuous:} $\delta$, $\mu$ (2 fitted).
\textit{Discrete:} $m \propto \zch^2$, $y = \sqrt{2}g$,
$\SU(3)_F$ with $\Nf{=}\Nc{=}3$.} \\
\hline
\end{tabular}
\end{table}

\paragraph{Strengths and limitations.}
Scenario~(a) uses standard, well-established Seiberg dynamics.
The ISS metastable vacuum provides a natural SUSY-breaking
mechanism with calculable lifetime (bounce action
$S \sim 10^2$--$10^3$).
The limitation: the top Yukawa is an external coupling
(partial compositeness), not generated by confinement, breaking
the ``everything from one meson matrix'' paradigm.

\subsection{Scenario (b): speculative Pati--Salam extension}

Scenario~(b) requires a Pati--Salam gauge group ($\Nc = 4$,
$\Nf = 6$) that is NOT part of the core sBootstrap framework.
In this scenario:
\begin{enumerate}
  \item The $\mathcal{N}{=}2$ $\SU(\Nc)$ theory with $\Nf$
    hypermultiplets has superpotential~\cite{HMS:1997,Fayet:1976}
    $W = \sqrt{2}\,g\,Q^i\Phi\tilde{Q}_i + \sqrt{2}\,m^j_i
    Q^i\tilde{Q}_j$, where the Yukawa $\sqrt{2}\,g$ equals the
    gauge coupling by $\mathcal{N}{=}2$ SUSY---it is not free.
  \item Adding an adjoint mass
    $\Delta W = \mu_\Phi\,\Tr\Phi^2$ breaks
    $\mathcal{N}{=}2 \to \mathcal{N}{=}1$ and generates the
    quartic~\eqref{eq:quartic}.
  \item At the baryonic root~\cite{Csaki:2011xn}, the magnetic
    adjoint VEV breaks the gauge group to
    $\SU(\Nf - \Nc) = \SU(2)$, which IS $\SU(2)_L$.
\end{enumerate}
The strength is that the Yukawa hierarchy is a hierarchy of VEVs
(Seiberg seesaw), not couplings.
The limitation: the Pati--Salam embedding requires gauge structure
beyond $\SU(3)_F$, the KK monopole identification for the top is
speculative, and the K\"ahler metric is still needed.
The results of Secs.~\ref{sec:composites}--\ref{sec:SO32}
are independent of Scenario~(b).
By Fradkin--Shenker complementarity, the ``UV $\mathcal{N}{=}2$''
picture (Higgs phase) and the ``emergent $\mathcal{N}{=}2$''
picture (confined phase) may describe the same physics, smoothly
connected.


%======================================================================
\section{Discussion}
\label{sec:discussion}
%======================================================================

\paragraph{M-theory lift.}
In the Hanany--Witten brane
construction~\cite{HananyWitten:1997}, $\Nc$ D4-branes are
suspended between two NS5-branes, with $\Nf$ D6-branes providing
quarks.
Mesons $M = Q\tilde{Q}$ are open D6--D6 strings; baryons
$B = \varepsilon QQQ$ are wrapped D4-branes.
Lifting to M-theory (the $x^{10}$ circle opens up):
D4-branes become M5-branes, D6-branes become Taub-NUT centres,
NS5-branes remain M5-branes.
The gauge theory lives on the M5 worldvolume wrapping a
holomorphic curve $\Sigma$ in the $(x^4, x^5, x^6, x^{10})$
space.
The meson (open D6--D6 string) lifts to an M2-brane ending on
the M5 worldvolume, coupling to $C_3$ (Dirac mass type).
The baryon (wrapped D4) lifts to an M5-brane wrapping $S^4$
with three string endpoints, coupling to $C_6$ (Majorana mass
type).

\begin{table}
\centering\small
\caption{M-theory dictionary for $\SU(3)$ SQCD composites.}
\label{tab:Mtheory}
\begin{tabular}{lllll}
\hline
Field theory & IIA & M-theory & $C$ & Mass \\
\hline
Meson $Q\tilde{Q}$ & D6--D6 & M2 & $C_3$ & Dirac \\
Baryon $\varepsilon Q^3$ & D4 & M5 & $C_6$ & Majorana \\
$\adj(M)$ & thr.\ D4 & thr.\ M5 & mixed & branch \\
Involution & T-dual & EM dual & $C_3{\leftrightarrow}C_6$ & --- \\
\hline
\end{tabular}
\end{table}
The involution $M_i \to \Lambda^4/M_i$ exchanges the A-period
($\langle M_i\rangle \propto 1/m_i$) with the B-period
($\adj(M)_i \propto m_i$) of the SW curve.
In M-theory this is electromagnetic duality: $C_3 \leftrightarrow
C_6$, exchanging M2 and M5 branes.
The direct/inverse mass assignment of Sec.~\ref{sec:seiberg} IS
the A/B duality.
The genus of the M5 worldvolume curve $\Sigma$ for $\SU(3)$ is
2, matching the two Cartan moduli $(z_1, z_2)$.

The M5-brane wrapping $\Sigma$ gives $\mathcal{N}{=}2$ in 4d:
the curve encodes the Seiberg--Witten prepotential.
The adjoint mass $\mu_\Phi$ deforms $\Sigma$, breaking
$\mathcal{N}{=}2 \to \mathcal{N}{=}1$.
At the deformed curve, the two chiral multiplets $(M_i,
\adj(M)_i)$ emerge as the two periods of $\Sigma$---the
geometric origin of the emergent hypermultiplet structure
of Sec.~\ref{sec:N2}.

The membrane dimensionality provides a consistency check: for a
$p$-brane whose worldvolume dimensions all scale with $\zch_i$,
$m_i \propto \zch_i^p$.
The ratio $K_p = \sum\zch_i^{2p}/(\sum\zch_i^p)^2$ is
direction-independent only for $p = 2$ (the M2-brane), by
Theorem~\ref{thm:n2}.
For $p = 5$ (M5-brane), $K_5$ depends on the Cartan angle.
This confirms the composite mesons have the right dimensionality
for M2-branes---not an independent derivation, since assuming
both worldvolume directions scale with $\zch_i$ is equivalent to
the bilinear structure.
Under $\SU(3)_c$, the 11d bosonic fields decompose as:
$C_3 \in \Lambda^3(9) = \mathbf{84}$ contains
$\mathbf{1}_C$, $\mathbf{3}_C$, $\overline{\mathbf{3}}_C$
but \textbf{never} $\mathbf{6}_C$; while
$g_{MN} \in \mathrm{Sym}^2(9) = \mathbf{45}$ contains
$\mathbf{6}_C$ as a gravitational modulus.
This explains the colour-sextet obstruction: $\mathbf{6}_C$
diquarks are in the gravity sector (Planck scale), not the
gauge sector (confinement scale).

The number $84 = \binom{9}{3}$ appears in five contexts:
(i)~$C_3$ polarisations ($\Lambda^3(\mathbb{R}^9)$);
(ii)~$E_8 \to \SU(9)$: $\mathbf{248} = \mathbf{80} + \mathbf{84} + \overline{\mathbf{84}}$;
(iii)~SM fermions: $96 = 84 + 12$;
(iv)~number theory: $|\text{disc}(\mathbb{Q}(\sqrt{-21}))| = 84$;
(v)~the binomial identity $\binom{9}{3} = 4 \times 21$.

\paragraph{Composite $\SU(2)_L$ (speculative).}
The CST construction~\cite{Csaki:2011xn} offers a mechanism in
which $\SU(2)_L$ is a residual magnetic gauge group at the
baryonic root, but requires a Pati--Salam embedding
($\Nc = 4$, $\Nf = 6$) beyond the core $\SU(3)_F$ framework.
An honest caveat: the Seiberg dual of $\SU(3)$ at $\Nf = 3$
has \emph{trivial} gauge group---there is no residual $\SU(2)$
to identify with $\SU(2)_L$.
The Pati--Salam embedding is a separate gauge sector that must
be grafted onto the framework; it does not arise from
$\SU(3)_F$ alone.
In the core sBootstrap framework, $\SU(2)_L$ is an
\textbf{elementary} gauge symmetry already present in
$G_{\rm UV}$ (eq.~\eqref{eq:GUV}), embedded via the branching
$\mathbf{3}_L \to \mathbf{2}_{-1/2} \oplus \mathbf{1}_{+1}$.
The CST construction offers a speculative \emph{alternative}
that is NOT required by the results of this paper.

\paragraph{Yukawa emergence.}
If the $\mathcal{N}{=}2$ relation $y = \sqrt{2}\,g$ is adopted
(Section~\ref{sec:N2}), the same gauge coupling produces both
branches: the F-term gives mesino masses $\propto \zch_i^2$
(direct), while the meson VEV Yukawa gives masses
$\propto 1/\zch_i^2$ (inverse).
In the holomorphic sector, flavour parameters are controlled by
$g$ (overall scale) and $\alpha$ (mass ratios); the Yukawa
hierarchy is not a hierarchy of couplings but of VEVs (Seiberg
seesaw).
The mapping to physical observables also requires the K\"ahler
metric $Z_i$, the adjoint mass $\mu_\Phi$, and the overall scale
$\mu \approx 1883\MeV$ (fitted, not derived).

\paragraph{Holomorphic vs.\ physical masses.}
The meson masses from the Seiberg vacuum are holomorphic: they
arise from the superpotential $W(M)$ and are protected by
holomorphy.
The physical fermion masses are $m_{\rm phys} = m_{\rm hol}/\sqrt{Z_i}$
where $Z_i$ is the K\"ahler metric.
For leptons, $Z$ is generation-universal (to $0.001\%$), so the
holomorphic masses directly determine the mass \emph{ratios}.
For quarks, $m_c/m_s$ does NOT run ($\gamma_m$ is
flavour-independent to 5-loop), and the lattice value
$m_c/m_s = 11.783 \pm 0.025$ differs from the lepton-derived
prediction $(2+\sqrt{3})^2 = 13.93$ by $18\%$.
This discrepancy dissolves when sectors are allowed different
Cartan angles: the chain seed gives
$\zch_s/\zch_c = 0.271$, within $1.1\%$ of the algebraic
ratio $r = 2 - \sqrt{3}$, without invoking K\"ahler
corrections (see ``Resolved'' in
Section~\ref{sec:open_problems}).

% REVISED: REF-001 — Kahler sensitivity discussion
\paragraph{K\"ahler sensitivity of quark-sector comparisons.}
All quark-sector numerical comparisons in
Table~\ref{tab:results} operate at the holomorphic level.
The physical masses are $m_{\rm phys} = m_{\rm hol}/\sqrt{Z_i}$,
where $Z_i$ is the unknown K\"ahler metric.
For leptons, $Z_i$ is generation-universal to $0.001\%$---the
precision of $K(e,\mu,\tau) = 2/3$---protected by the emergent
$\mathcal{N}{=}2$ structure (Sec.~\ref{sec:open_problems}).
For quarks, which enter through partial
compositeness~\cite{Kaplan:1991dc}
(eq.~\ref{eq:mq_partial}), the physical mass is
$m_q \sim \lambda_{q_L}\lambda_{q_R}\langle\mathcal{O}\rangle/M_*$,
where $\lambda_{q_L,R}$ are elementary--composite mixing
parameters.
The $\mathcal{N}{=}2$ protection does not obviously extend to
partial compositeness masses, so $Z_i$ may receive
$\mathcal{O}(1$--$10\%)$ generation-dependent corrections.
These could shift $K^{-1}(d,s,b)$ by up to $\sim\!1\%$, the bion
mass $m_b$ by up to $\sim\!5\%$, and $m_t$ by up to
$\sim\!10\%$.
The cross-sector product
rule~\eqref{eq:cross_sector}---already violated by a factor 4.6
at the holomorphic level---demonstrates that
generation-dependent K\"ahler corrections across the
lepton--quark boundary are large.
\emph{Until the K\"ahler metric is computed or bounded, the
quark-sector holomorphic agreements should be regarded as
suggestive, not as precision tests of the framework.}

\paragraph{Why $m_i = \zch_i^2$: honest status.}
The charge-squared law is motivated by three structural inputs:
holomorphy ($W(M)$ holomorphic $\Rightarrow$ mass
$\propto \zch^2$), bilinear structure ($M = Q\tilde{Q}
\Rightarrow$ two preon factors), and canonical normalisation
(Theorem~\ref{thm:n2}).
However, seven attempts to derive $m = \zch^2$ from SUSY-breaking
dynamics all fail for the following structural reason: gauge
mediation gives $\zch^2$ scaling for \emph{soft} (K\"ahler)
masses, but we need it for \emph{holomorphic} (superpotential)
masses---the protected masses that determine exact $K = 2/3$.
The diagonal mesons $M^i_i$ are $\UU(1)_F$-neutral, which kills
D-term approaches; and soft masses enter the scalar potential,
not the fermion mass matrix.
The charge-squared law remains the best-motivated ansatz, not a
theorem of the confining theory.
The Coulomb self-energy argument~\cite{Koide:1983qe} (a capacitor
energy $E = Q^2/(2C)$) is the physically most transparent route:
it is rigorous for weakly-coupled Abelian fields but must be
assumed to survive in the strongly-coupled confining regime.

\paragraph{Hierarchy reframing.}
In the sBootstrap, SUSY is broken at $f = f_\pi \sim 92\MeV$
---the QCD scale, not the TeV scale.
The electroweak VEV $v_{\rm EW} = 246\GeV$ is a \emph{derived}
quantity: $v_{\rm EW} \sim m_t / y_t \sim 173/0.7 \sim 247\GeV$.
The apparent hierarchy $v_{\rm EW}/\Lambda_{\rm QCD} \sim 10^3$
is the Yukawa hierarchy---which is the flavour problem that
the charge-squared law addresses.
This is a \emph{reframing}, not a solution: the hierarchy is
moved from $v_{\rm EW}/M_{\rm Pl}$ to $y_t/y_e$, and the latter
still requires explanation (the Cartan angle $\alpha$, which is
not derived).

\paragraph{Strong CP and the Nelson--Barr mechanism.}
Identifying the involution $M_i \to \Lambda^4/M_i$ with CP
gives a Nelson--Barr~\cite{Nelson:1984,Barr:1984} solution to
the strong CP problem.
CP is an exact symmetry of the UV theory (the involution is a
symmetry of the quantum constraint); it is spontaneously broken
by the meson VEVs $\langle M_i\rangle$ (which are real but
unequal).
At tree level: $\arg\det m_q = 0$ (the superpotential is real
under the involution).
Radiative corrections generate
$\delta\bar\theta \sim (\alpha/4\pi)^2 J \sim (10^{-3})^2
\times 3 \times 10^{-5} \sim 10^{-11}$,
nine orders of magnitude below the experimental bound
$|\bar\theta| < 10^{-10}$.

\paragraph{The AM-GM scale chain.}
The framework produces three mass scales from $\SU(3)$
confinement and the bilinear structure $M = Q\tilde{Q}$.
Define the arithmetic and geometric means of the lepton charges:
$\mathrm{AM} = \frac{1}{3}\sum_i \zch_i$ and
$\mathrm{GM} = (\prod_i \zch_i)^{1/3}$.
Then:
\begin{align}
  M_0 &= \mathrm{AM}^2 = 313.8\MeV
    && (\text{Koide scale} \approx m_p/3)\,,\\
  \tfrac{1}{2}f_\pi &= \mathrm{GM}^2 = 45.8\MeV
    && (\text{SUSY breaking})\,,\\
  \mu &= 6\,M_0 = 1883\MeV
    && (\text{mass scale} \approx 2\,m_p)\,.
\end{align}
The $f_\pi$ relation $f_\pi = 2(\zch_e\,\zch_\mu\,\zch_\tau)^{2/3}$
holds to $0.7\%$.
The ratio $M_0/f_\pi = \frac{1}{2}(\mathrm{AM}/\mathrm{GM})^2$
measures the hierarchy: $\mathrm{AM}/\mathrm{GM} = 2.618$ encodes
the $\tau/e$ mass ratio.
The sum $\mu = m_e + m_\mu + m_\tau = 1883\MeV$ coincides with
$2\,m_p = 1877\MeV$ ($0.3\%$).
Both arise from $\SU(3)$ confinement at similar scales; the
probability of some match at $0.3\%$ among $\sim\!50$ hadron
masses is $p \approx 15\%$ after look-elsewhere correction.
We note the coincidence without claiming it is derived.


%======================================================================
\section{Open problems and conclusions}
\label{sec:open}
%======================================================================

\subsection{Numerical results}

% REVISED: REF-001 — Kahler sensitivity column added
\begin{table*}
\centering\small
\begin{tabular}{llllll}
\hline
Observable & Method & Predicted & PDG 2024 & Deviation
  & K\"ahler sensitivity \\
\hline
$K(e,\mu,\tau)$ & trace id.\ & $2/3$ & $0.666661$
  & $0.001\%$ & protected ($\delta Z/Z \sim 10^{-9}$) \\
$K(-\pi,D_s,B)$ & meson & $\approx 2/3$ & $0.6674$
  & $0.10\%$ & protected (HQET: $Z$ cancels in ratio) \\
$K(-\pi,D,B)$ & meson & $\approx 2/3$ & $0.673$
  & $0.91\%$ & protected (HQET) \\
$K^{-1}(d,s,b)$ & seesaw & $\approx 2/3$ & $0.6652$
  & $0.22\%$ & \textbf{uncontrolled}
  ($\delta K \sim 1$--$10\%$ from $\delta Z_i$) \\
$m_b$ (bion) & eq.~\eqref{eq:v0_doubling_intro} & $4183\MeV$ & $4183$
  & $0.01\%$ & \textbf{uncontrolled}
  ($\delta m_b/m_b \sim \delta\sqrt{Z_b}$) \\
$m_t$ (compl.) & eq.~\eqref{eq:top_completion} & $168.5\GeV$ & $172.4$
  & $2.3\%$ & \textbf{uncontrolled}
  ($\delta m_t \sim 5$--$10\%$ possible) \\
$|V_{us}|$ & $\sqrt{m_\mu/m_\tau}$ & $0.244$ & $0.224$
  & $8.7\%$ & cross-sector: needs $Z_d/Z_s \approx 1.18$ \\
\hline
\end{tabular}
\caption{Quantitative results with K\"ahler sensitivity.
  All entries use PDG 2024~\cite{PDG:2024} inputs.
  ``Protected'' means the comparison involves ratios within a sector
  where $\mathcal{N}{=}2$ K\"ahler protection applies
  (Sec.~\ref{sec:open_problems}).
  ``Uncontrolled'' means the holomorphic-to-physical mapping
  depends on unknown generation-dependent K\"ahler corrections
  from partial compositeness; the stated deviation is holomorphic
  and may shift by $\mathcal{O}(Z_i)$.}
\label{tab:results}
\end{table*}

\paragraph{Tier~1: high precision ($< 0.1\%$).}
$K(e,\mu,\tau) = 0.666661$ ($0.001\%$) stands alone as the
most precise mass relation in particle physics outside the
gauge sector.
The bion prediction $\zch_b = 3\zch_s + \zch_c$ ($0.01\%$) is
also in this tier.
These are the two results for which accidental coincidence is
least plausible: a joint Monte Carlo with $10^7$ random mass
triples gives $p < 10^{-7}$.

\paragraph{Tier~2: moderate precision ($0.1$--$1\%$).}
The meson triple $K(-\pi, D_s, B) = 0.6674$ ($0.10\%$) and the
inverse triple $K^{-1}(d,s,b) = 0.6652$ ($0.22\%$).
The cross-sector product rule predicts
$|V_{us}|_{\rm hol} = \sqrt{m_\mu/m_\tau} = 0.244$ ($8.7\%$ off),
which is an order-of-magnitude test rather than a precision result.
At $Q = 100\TeV$, the inverse triple $K^{-1}(s,c,-t) = 0.6666$
($0.004\%$) joins Tier~1 under RG running.

\paragraph{Tier~3: the chain as a correlated set.}
The quark chain $(t,b,c)$--$(b,c,-s)$--$(c,s,0)$--$(s,0,u{+}d)$
is individually less precise than the lepton Koide ($0.3$--$1.2\%$
per link).
But the chain is not four independent observations: it is a
\textbf{correlated set of overlapping triples} sharing members.
The four links use only 6 quarks and contain
$\binom{6}{3} = 20$ possible triples, of which exactly 4
(the consecutive windows) have $K \approx 2/3$ simultaneously.
% REVISED: REF-003 — Monte Carlo methodology documented
The probability that a random mass spectrum produces four
consecutive windows all within $1.2\%$ of $2/3$ is
$p \sim 1/21\,000$, estimated as follows.
We draw six quark masses independently from a log-uniform
distribution over $[m_u, m_t] = [2.2\MeV, 172.4\GeV]$, assign
random signs ($\pm 1$, uniform) to each mass in the Koide formula,
and evaluate $K$ on the four consecutive windows of three.
A ``hit'' requires all four windows to satisfy
$|K - 2/3|/(2/3) < 0.012$ simultaneously.
Over $10^6$ trials, we find $47 \pm 7$ hits, giving
$p = (4.7 \pm 0.7) \times 10^{-5} \approx 1/21\,000$.
This is \emph{before} any look-elsewhere correction; the chain
structure (consecutive windows with a fixed ordering) is
specified a priori, not selected from the data.
The chain is a single coherent observation, and its significance
should be assessed jointly, not link by link.

\paragraph{The inverse chain at high scale.}
Under 5-loop SMDR running, two inverse triples converge on
$K^{-1} = 2/3$ at $100\TeV$:
$K^{-1}(d,s,b) = 0.6667$ ($0.012\%$) and
$K^{-1}(s,c,-t) = 0.6666$ ($0.004\%$).
Both share the strange quark as pivot, forming an inverse chain
$(d,s,b)$--$(s,c,-t)$ shorter than the direct chain but more
precise at high scales.

\paragraph{Prediction vs.\ postdiction.}
The lepton Koide (Koide 1983~\cite{Koide:1983qe}) and the
inverse quark Koide (Rivero 2005~\cite{Rivero:2005b}) were
observed before the present framework; they are postdictions.
The bion relation $\zch_b = 3\zch_s + \zch_c$ is also a
postdiction (identified after seeing data).
The framework's genuine predictions---testable going forward---are
the LFV ratio and the family gauge boson masses.
The CKM matrix beyond the Cabibbo angle is not currently
predicted.


\subsection{Open problems}
\label{sec:open_problems}

\paragraph{Resolved or partially resolved.}
\begin{itemize}
\item \textbf{The fermion K\"ahler metric} (formerly ``master
  obstruction'').
  The emergent $\mathcal{N}{=}2$ hypermultiplet structure
  protects the fermion K\"ahler metric: the $\mathcal{N}{=}2$
  breaking correction is $\delta Z/Z \sim g^2\Lambda^2/\mu_\Phi^2
  \sim 5 \times 10^{-9}$, negligible.
  This \emph{explains} the $0.001\%$ lepton Koide: it is
  $\mathcal{N}{=}2$ protection, not a mystery.
  The fermion $Z_i$ is generation-universal.
  \textbf{This part of Problem~1 is solved.}

\item \textbf{The $m_c/m_s$ ratio} (formerly part of Problem~1).
  The holomorphic ratio $m_c/m_s = (2+\sqrt{3})^2 = 13.93$
  vs.\ lattice $11.78$ ($18\%$) was attributed to
  $Z_c/Z_s \approx 1.37$.
  However, $m_c/m_s$ does NOT run ($\gamma_m$
  flavour-independent to 5-loop; $< 1\%$ drift from
  $2\GeV$ to $282\TeV$).
  The ``discrepancy'' dissolves when sectors are allowed
  different Cartan angles: the chain seed gives
  $\zch_s/\zch_c = 0.271$, within $1.1\%$ of the algebraic
  ratio $r = 2 - \sqrt{3}$.
  No K\"ahler correction is needed.

\item \textbf{The $Q_2$ operator} (formerly ``open: can the
  $\mathbb{Z}_2$ be promoted?'').
  The operator $R_{ij} = \langle M_k\rangle$
  ($k \neq i,j$; $R_{ii} = 0$) is explicitly constructed.
  It satisfies the \emph{exact} identity
  \begin{equation}
    R \cdot \mathbf{M} = 2\,\adj(\mathbf{M})
    \label{eq:RM_adjM}
  \end{equation}
  (verified symbolically): $R$ acting on the meson VEV vector
  gives twice the adjugate vector.
  Its Cayley--Hamilton identity on the constraint surface is
  \begin{equation}
    R^3 = \Tr(D^2)\,R + 2\Lambda^6\,I\,,
    \label{eq:R_cubic}
  \end{equation}
  where $D = \diag(\langle M_i\rangle)$.
  This is a \emph{cubic} minimal polynomial, not the quadratic
  $R^2 = c\,I$ of standard $\mathcal{N}{=}2$.
  The normalised operator $T = R/\sqrt{\det D}$ has eigenvalues
  $\approx \{+1.07, -3 \times 10^{-6}, -1.07\}$ for the
  lepton hierarchy---an ``almost-involution'' with one
  near-zero eigenvalue.
  Whether the cubic algebra~\eqref{eq:R_cubic} defines a
  \emph{deformed} $\mathcal{N}{=}2$ structure (controlled by
  the mass hierarchy $\Tr D^2$) is an open question.
\end{itemize}

\paragraph{Remaining open problems.}
\begin{enumerate}
\item \textbf{The Cartan angle $\alpha \approx 12.7^\circ$}
  (absorbs old Problems 3 and 4).
  One-loop CW analysis shows all perturbative contributions to
  $V(\delta)$ minimise at $\delta = 120^\circ$ (the democratic
  point), not at $\delta_{\rm phys} = 132.73^\circ$: the
  tree-level quartic gives
  $V \propto \mathrm{Var}(1/\zch_i^2)$, and the CW correction
  has the same angular profile
  ($\mathrm{STr}\,\mathcal{M}^4 \propto V_{\rm quartic}$).
  The deviation $\alpha = 12.73^\circ$ from the democratic
  point is a flat direction of the perturbative theory.
  Its value requires either the non-perturbative K\"ahler
  metric for the \emph{scalar} meson VEV (not the fermion
  metric, which is solved) or an as-yet-unknown selection
  mechanism.

\item \textbf{The CKM matrix.}
  The off-diagonal meson mass matrix block-decouples from the
  diagonal mesons at tree level (the $10 \times 10$ fermion
  mass matrix decomposes into a $4 \times 4$ diagonal block
  and three $2 \times 2$ off-diagonal blocks).
  CKM mixing therefore \emph{cannot} arise from tree-level
  meson mixing; it requires K\"ahler corrections, loop effects,
  or RG running.
  The holomorphic cross-sector prediction
  $|V_{us}| = \sqrt{m_\mu/m_\tau} = 0.244$ is $8.7\%$ off,
  suggesting K\"ahler corrections $Z_d/Z_s \approx 1.2$.

\item \textbf{PMNS and neutrino masses.}
  The baryons $B, \tilde{B}$ are SM singlets (sterile fermions,
  Table~\ref{tab:IR_composites}).
  A tree-level seesaw gives $M_\nu \propto m_\ell^2/M_B$, but
  the tree-level PMNS matrix is diagonal---contradicting the
  observed large mixing.

\item \textbf{The cubic $Q_2$ algebra.}
  Eq.~\eqref{eq:R_cubic} shows $R$ satisfies a cubic, not a
  quadratic, minimal polynomial.
  The coefficient $2\Lambda^6 = 2\det M$ is the quantum
  constraint; $\Tr(D^2) = \sum M_i^2$ is the quadratic
  Casimir.
  Does this cubic define a consistent deformed
  $\mathcal{N}{=}2$ superalgebra?
  If so, the ``emergent $\mathcal{N}{=}2$'' label becomes
  precise: a cubic deformation of the hypermultiplet algebra
  controlled by the family hierarchy.

\item \textbf{The $\SO(32)$ projection.}
  The $\mathbf{496}$ contains 442 states beyond the sBootstrap
  content; the projection mechanism is unknown.
\end{enumerate}


\subsection{Falsification conditions}

\begin{itemize}
  \item The LFV ratio
    $\mathrm{Br}(\tau \to e\gamma)/\mathrm{Br}(\tau \to \mu\gamma)
    = (\zch_e/\zch_\mu)^4 \approx 5.5 \times 10^{-10}$ is a
    parameter-free prediction testable at Belle~II.
  \item Family gauge bosons at $10^2$--$10^3\TeV$ produce
    flavour-changing effects testable at FCC-hh.
  \item Lattice K\"ahler metric strongly generation-dependent
    for leptons ($Z_\mu/Z_e \neq 1$ at $1\%$) would kill the
    framework.
  \item LFV ratio deviating from $(\zch_e/\zch_\mu)^4$ by more
    than a factor~2 falsifies the $\UU(3)_F$ charge assignment.
  \item Discovery of a fourth generation contradicts $N = 3$
    uniqueness.
\end{itemize}

\subsection{Conclusions}

We have presented a framework in which Standard Model fermions
are identified with composite fermions of a confining $\SU(3)_F$
family gauge theory: mesinos are leptons, and the adjugate
composites provide quark partners through partial compositeness.
The confined spectrum automatically produces a pair
$(M_i, \adj(M)_i)$ per family with the mathematical structure of
an $\mathcal{N}{=}2$ hypermultiplet---emergent from confinement,
not imposed in the UV.

The framework rests on one central ansatz---the charge-squared
law $m_i = \mu\,\zch_i^2$---which is the unique monomial
compatible with direction-independent $K = 2/3$
(Theorem~\ref{thm:n2}), supported by the Coulomb self-energy
of the composite~\cite{Koide:1983qe} and by the bilinear
structure $M = Q\tilde{Q}$.
We have been honest that this is the best-motivated ansatz, not
a theorem of the confining theory: seven attempts to derive it
from SUSY-breaking dynamics all fail for holomorphic masses.

What is \textbf{established}:
\begin{enumerate}
  \item \textbf{Composites:} 9 mesinos furnish $3L_i + 3e_i^c$
    (exact group theory).
    Both Higgs doublets from one meson matrix:
    $H_d = M^{1,2}$, $H_u = \adj(M)^{1,2}$, ratio fixed by the
    Seiberg constraint.
    The ``superpartners'' are the mesons, split from mesinos by
    $f_\pi$.

  \item \textbf{The adjugate mechanism:} $M \propto 1/m$ and
    $\adj(M) \propto m$ provide two Yukawa operators from one
    composite.
    A single UV condition produces both standard lepton and
    inverse quark $K = 2/3$.

  \item \textbf{Energy-balance observations:}
    $K(e,\mu,\tau) = 0.666661$ ($0.001\%$);
    $K(-\pi, D_s, B) = 0.6674$ ($0.10\%$);
    $K^{-1}(d,s,b) = 0.6652$ ($0.22\%$);
    and at $Q = 100\TeV$, $K^{-1}(s,c,-t) = 0.6666$ ($0.004\%$).
    The quark chain
    $(t,b,c)$--$(b,c,-s)$--$(c,s,0)$--$(s,0,u{+}d)$
    has $K$ within $1.2\%$ in every link (joint $p \sim
    1/21\,000$), and is compressed by three algebraic relations
    (seed map, iterated seed, bion/threshold) plus a
    positive-completion formula that predicts
    $m_t = 168.5\GeV$ ($2.3\%$).

  \item \textbf{Three mass scales from one confinement:}
    $M_0 = \mathrm{AM}^2 = 313.8\MeV \approx m_p/3$;
    $f_\pi/2 = \mathrm{GM}^2$;
    $\mu = 6\,M_0 = 1883\MeV \approx 2\,m_p$.
    The Fermi scale and the QCD scale are \emph{not} independently
    set---both arise from $\SU(3)$ confinement.

  \item \textbf{Generation uniqueness:} five criteria
    (two structural---Diophantine and Casimir--Koide---and three
    suggestive) select $N = 3$; composite content selects
    $\SO(32)$.
\end{enumerate}

The seven open problems of Section~\ref{sec:open_problems}
constitute what the framework does \textbf{not yet} provide;
the most pressing are the Cartan angle (Problem~1), the
CKM matrix (Problem~2), and the cubic $Q_2$ algebra
(Problem~4).

On the experimental side, the LFV ratio
$\mathrm{Br}(\tau \to e\gamma)/\mathrm{Br}(\tau \to \mu\gamma)
= (\zch_e/\zch_\mu)^4$ is a parameter-free prediction testable
at Belle~II, and family gauge bosons at $10^2$--$10^3\TeV$
produce flavour-changing effects accessible to a future
$100\TeV$ collider.

\paragraph{Adjugate crossing.}
Under RG running, the charged-lepton energy-balance
$K(e,\mu,\tau)$ and the inverse down-quark energy-balance
$K^{-1}(d,s,b)$ both remain close to $2/3$ but at slightly
different values.
Remarkably, they cross at $Q \approx 15.7\GeV$: at this scale,
$K_{\rm lep} = K^{-1}_{\rm down} = 0.6663$, and the gap at
$m_b$ is only $0.06\%$.
This crossing is a quantitative consequence of the adjugate
mechanism: both $K$ and $K^{-1}$ are determined by the same
charges, and their coincidence at a particular RG scale reflects
the interplay of direct and inverse operators under running.
The crossing scale $Q \approx 15.7\GeV$ is near the $b$-quark
threshold, consistent with the phase transition from $\Nf = 4$
to $\Nf = 3$ in Table~\ref{tab:phases}.

The framework represents a structural EFT, not a UV completion:
it organises the flavour sector using composites of a confining
gauge theory, with exact results (group theory, holomorphy,
anomaly matching) supplemented by a well-motivated ansatz
(the charge-squared law) and empirically verified patterns
(the chain, the bion relation, the seed structure).
The gap between what is established and what is not is clearly
delineated by the seven open problems listed above.


%======================================================================
\begin{acknowledgments}
The author thanks M.~Porter for sustained theoretical discussions
over 2011--2024 that identified Seiberg duality and the mesino
interpretation as the mechanism underlying the sBootstrap,
and C.~Brannen for the circulant mass matrix parametrisation.
During the preparation of this work the author used Claude
(Anthropic) for algebraic verification, numerical checks,
and drafting assistance.
The author reviewed and edited all content and takes full
responsibility for the publication.
\end{acknowledgments}

% REVISED: REF-004 — Note Added trimmed; CKM/PMNS details deferred
% to companion letter
\paragraph*{Note added.}
\emph{(i)~Cayley--Hamilton.}
The cubic algebra $R^3 = \mathrm{Tr}(D^2)\,R + 2\Lambda^6\,I$ of
Sec.~\ref{sec:open_problems} coincides with the Cayley--Hamilton
equation of the $\mathcal{N}{=}2$ adjoint scalar~$\Phi$ under
$u_2 = \mathrm{Tr}(D^2)$, $u_3 = 2\det(D)$.
\emph{(ii)~Seiberg--Witten curve.}
The SW curve for $\SU(3)$ with $\Nf=3$ at the Koide point
is a smooth genus-2 surface; in the democratic limit it develops a
monopole singularity at $\Lambda = m = s_1/3$.
The Koide constraint reduces the curve to a one-parameter family
in the Cartan angle~$\delta$.
\emph{(iii)~Cabibbo angle.}
Koide's original Cabibbo prediction~\cite{Koide:1983qe}
$\tan\theta_C = \sqrt{3}(\sqrt{m_\mu} - \sqrt{m_e})/
(2\sqrt{m_\tau} - \sqrt{m_\mu} - \sqrt{m_e})$
is algebraically identical to the Cartan angle~$\delta$
($\delta = 0.222$, $1.2\%$ from PDG).
This extends the quark--lepton complementarity of
Raidal~\cite{Raidal:2004} and the quark Koide phases of
Zenczykowski~\cite{Zenczykowski:2012}.
The full CKM and PMNS parametrisations from $\delta$---including
$|V_{us}| = 2/9$, $|V_{cb}| = 4/99$, and PMNS deviations from
tribimaximal mixing---are phenomenological observations
(not yet derived from the $\SU(3)_F$ dynamics) and are presented
in a companion letter~\cite{companion:PRL}.
\emph{(iv)~RG comparison.}
Using multi-loop SM running~\cite{AHS:2024,MartinRobertson:2019},
$K^{-1}(d,s,b) = 2/3$ at $Q \approx 10\TeV$;
in the MSSM, $K^{-1}$ at the GUT scale is $0.23$--$0.65\%$
above $2/3$ for all $\tan\beta$ and $M_{\rm SUSY}$, disfavouring
conventional supersymmetric extensions.


%======================================================================
\begin{thebibliography}{99}

\bibitem{Chew:1962}
  G.F.~Chew,
  Rev.\ Mod.\ Phys.\ \textbf{34}, 394 (1962).

\bibitem{Miyazawa:1966}
  H.~Miyazawa,
  Prog.\ Theor.\ Phys.\ \textbf{36}, 1266 (1966).

\bibitem{CattoGursey:1985}
  S.~Catto and F.~G\"ursey,
  Nuovo Cim.\ A \textbf{86}, 201 (1985).

\bibitem{Rivero:2024epjc}
  A.~Rivero,
  Eur.\ Phys.\ J.\ C \textbf{84}, 1058 (2024),
  arXiv:2407.05397.

\bibitem{VolkovAkulov:1973}
  D.V.~Volkov and V.P.~Akulov,
  Phys.\ Lett.\ B \textbf{46}, 109 (1973).

\bibitem{FradkinShenker:1979}
  E.~Fradkin and S.H.~Shenker,
  Phys.\ Rev.\ D \textbf{19}, 3682 (1979).

\bibitem{Kaplan:1991dc}
  D.B.~Kaplan,
  Nucl.\ Phys.\ B \textbf{365}, 259 (1991).

\bibitem{NelsonStrassler:2000}
  A.E.~Nelson and M.J.~Strassler,
  JHEP \textbf{09}, 030 (2000),
  arXiv:hep-ph/0006251.

\bibitem{Seiberg:1994bz}
  N.~Seiberg,
  Nucl.\ Phys.\ B \textbf{435}, 129 (1995),
  arXiv:hep-ph/9411149.

\bibitem{CsakiSchmaltzSkiba:1997}
  C.~Cs\'aki, M.~Schmaltz, and W.~Skiba,
  Phys.\ Rev.\ Lett.\ \textbf{78}, 799 (1997),
  arXiv:hep-th/9610139.

\bibitem{Tong:SUSY}
  D.~Tong,
  \textit{Supersymmetric Field Theory},
  lecture notes, University of Cambridge (2022),
  \url{http://www.damtp.cam.ac.uk/user/tong/susy.html}.

\bibitem{ISS:2006}
  K.~Intriligator, N.~Seiberg, and D.~Shih,
  JHEP \textbf{04}, 021 (2006),
  arXiv:hep-th/0602239.

\bibitem{SeibergWitten:1994}
  N.~Seiberg and E.~Witten,
  Nucl.\ Phys.\ B \textbf{426}, 19 (1994),
  arXiv:hep-th/9407087.

\bibitem{HMS:1997}
  T.~Hirayama, N.~Maekawa, and S.~Sugimoto,
  arXiv:hep-th/9705069 (1997).

\bibitem{Csaki:2011xn}
  C.~Csaki, Y.~Shirman, and J.~Terning,
  Phys.\ Rev.\ D \textbf{84}, 095011 (2011),
  arXiv:1106.3074.

\bibitem{Koide:1983qe}
  Y.~Koide,
  Phys.\ Lett.\ B \textbf{120}, 161 (1983).

\bibitem{Koide:1990}
  Y.~Koide,
  Mod.\ Phys.\ Lett.\ A \textbf{5}, 2319 (1990).

\bibitem{GellMannOakesRenner:1968}
  M.~Gell-Mann, R.J.~Oakes, and B.~Renner,
  Phys.\ Rev.\ \textbf{175}, 2195 (1968).

\bibitem{Sumino:2009}
  Y.~Sumino,
  JHEP \textbf{05}, 075 (2009),
  arXiv:0812.2103.

\bibitem{PDG:2024}
  S.~Navas \textit{et al.} (Particle Data Group),
  Phys.\ Rev.\ D \textbf{110}, 030001 (2024).

\bibitem{MartinRobertson:2019}
  S.P.~Martin and D.G.~Robertson,
  Phys.\ Rev.\ D \textbf{100}, 073004 (2019),
  arXiv:1907.02500.

\bibitem{Rivero:2005b}
  A.~Rivero,
  arXiv:hep-ph/0510068 (2005).

\bibitem{GattoSartoriTonin:1968}
  R.~Gatto, G.~Sartori, and M.~Tonin,
  Phys.\ Lett.\ B \textbf{28}, 128 (1968).

\bibitem{ORaifeartaigh:1975}
  L.~O'Raifeartaigh,
  Nucl.\ Phys.\ B \textbf{96}, 331 (1975).

\bibitem{Unsal:2008}
  M.~\"Unsal,
  Phys.\ Rev.\ Lett.\ \textbf{100}, 032005 (2008),
  arXiv:0708.1772.

\bibitem{Shih:2008}
  D.~Shih,
  JHEP \textbf{02}, 091 (2008),
  arXiv:hep-th/0703196.

\bibitem{HananyWitten:1997}
  A.~Hanany and E.~Witten,
  Nucl.\ Phys.\ B \textbf{492}, 152 (1997),
  arXiv:hep-th/9611230.

\bibitem{Nelson:1984}
  A.E.~Nelson,
  Phys.\ Lett.\ B \textbf{136}, 387 (1984).

\bibitem{Barr:1984}
  S.M.~Barr,
  Phys.\ Rev.\ Lett.\ \textbf{53}, 329 (1984).

\bibitem{Fayet:1976}
  P.~Fayet,
  Nucl.\ Phys.\ B \textbf{113}, 135 (1976).

\bibitem{MasieroVeneziano:1985}
  A.~Masiero and G.~Veneziano,
  Nucl.\ Phys.\ B \textbf{249}, 593 (1985).

\bibitem{LutyMohapatra:1997}
  M.A.~Luty and R.N.~Mohapatra,
  Phys.\ Lett.\ B \textbf{396}, 161 (1997),
  arXiv:hep-ph/9611343.

\bibitem{Kostant:1959}
  B.~Kostant,
  Amer.\ J.\ Math.\ \textbf{81}, 973 (1959).

% REVISED: REF-009
\bibitem{LutyTerning:2000}
  M.A.~Luty and J.~Terning,
  Phys.\ Rev.\ D \textbf{62}, 075006 (2000),
  arXiv:hep-ph/9812482.

\bibitem{Karasik:2020}
  A.~Karasik,
  JHEP \textbf{02}, 166 (2021),
  arXiv:2012.11154.

\bibitem{AHS:2024}
  S.~Antusch, C.~Hinze, and S.~Saad,
  arXiv:2510.01312 (2025).

\bibitem{companion:PRL}
  A.~Rivero,
  ``CKM and PMNS mixing from the Koide Cartan angle,''
  companion letter (2025).

\bibitem{companion:N2}
  A.~Rivero,
  ``Emergent $\mathcal{N}{=}2$ structure in the sBootstrap,''
  companion paper (2025).

\bibitem{Raidal:2004}
  M.~Raidal,
  Phys.\ Rev.\ Lett.\ \textbf{93}, 161801 (2004),
  arXiv:hep-ph/0404046.

\bibitem{Zenczykowski:2012}
  P.~Zenczykowski,
  Phys.\ Rev.\ D \textbf{86}, 117303 (2012),
  arXiv:1210.4125.

\end{thebibliography}

\end{document}
